\documentclass[12pt,a4paper]{article}
\usepackage[T1]{fontenc}
\usepackage[utf8]{inputenc}
\usepackage[ngerman,latin]{babel}
\usepackage{amsmath,amssymb,amsthm}
\usepackage{geometry}
\usepackage{titlesec}
\usepackage{fancyhdr}
\usepackage{longtable}
\usepackage{booktabs}
\usepackage{array}
\usepackage{multirow}
\usepackage{graphicx}
\usepackage{float}

\geometry{
  a4paper,
  left=2.5cm,
  right=2.5cm,
  top=2.5cm,
  bottom=2.5cm
}

\pagestyle{fancy}
\fancyhf{}
\fancyhead[C]{\small Carl Friedrich Gau\ss{} Werke — Band 5}
\fancyfoot[C]{\thepage}
\renewcommand{\headrulewidth}{0.4pt}

\titleformat{\section}{\Large\bfseries}{\thesection.}{0.5em}{}
\titleformat{\subsection}{\large\bfseries}{\thesubsection.}{0.5em}{}

\title{Carl Friedrich Gau\ss{} Werke --- Band 5, Teil 11}
\author{Carl Friedrich Gau\ss}
\date{}

\begin{document}

\maketitle
\thispagestyle{empty}

\tableofcontents
\newpage

%==============================================================================
% SECTION: Beobachtungen der magnetischen Inclination in Göttingen (forts.)
% Pages 501-504
%==============================================================================
\section*{Beobachtungen der magnetischen Inclination in Göttingen}
\addcontentsline{toc}{section}{Beobachtungen der magnetischen Inclination in Göttingen}

\subsection*{Fortsetzung}
\addcontentsline{toc}{subsection}{Fortsetzung}

\noindent
\ldots{} einzelnen Tage lassen sich zwar diese Bestandtheile nicht von einander scheiden, 
allein eine Abschätzung eines Mittelwerths der wirklichen Schwingungen mag 
bei einer so zahlreichen Reihe wohl versucht werden. In dieser Absicht habe ich 
zuvörderst die Inclinationen unter Voraussetzung einer regelmässigen jährlichen 
Abnahme von 3 Minuten auf den 21.\ Junius reducirt, und dann die Quadrate der 
Differenzen von dem Mittelwerthe addirt; diese Summe 220184 mit 30 dividirt 
gibt 7339,5 als Quadrat des mittlern Fehlers, dem man sich aussetzt, wenn man 
aufs Gerathewohl eine jener 31 Inclinationen als die mittlere für die Zeit der Be-
obachtung gültige ansehen wollte. Soll die ungleiche Zuverlässigkeit der drei 
Beobachtungsgruppen berücksichtigt werden, so ergeben die Grundsätze der Wahr-
scheinlichkeitsrechnung, indem man den mittlern Fehler für die vier ersten Be-
obachtungen mit $m'$, für die drei letzten mit $m''$, und für die 24 übrigen mit $m$, 
das mittlere Schwanken der Inclination selbst aber mit $M$ bezeichnet, folgende 
Gleichung:
\[
7339{,}5 = \tfrac{24mm + 4m'm' + 3m''m'' + 31MM}{31}
\]
Für $mm$ ist oben der Werth 1829,25 gefunden, oder es kann wenigstens diese 
Zahl wie eine hinlängliche Annäherung angesehen werden, für die sieben an-
dern Beobachtungen mag in Ermangelung eines sichern Maassstabes die Zahl der 
Einstellungen, woraus die Resultate abgeleitet sind, zum Grunde gelegt, also 
$m''m'' = mm$, $m'm' = 8mm$
gesetzt werden. Dadurch wird
\[
MM = \frac{7339{,}5 - 434\cdot 1829{,}25}{31} = 5168
\]
und $M = 71{,}9$.

\subsection*{32.}
\addcontentsline{toc}{subsection}{32.}

Mit demselben Instrumente und an demselben Platze hatte ich auch schon 
im vorigen Jahre eine Reihe von Inclinationsbeobachtungen gemacht, von denen 
ich jedoch nur die Endresultate hieher setze.

\begin{center}
\begin{tabular}{r@{ }l@{ }l@{ }l}
1841 & Sept.\ 22 & \ldots\ & $67^\circ\; 40'\; 20''$ \\
& 24 & \ldots\ & 40\quad 53 \\
& 26 & \ldots\ & 41\quad 41 \\
& 27 & \ldots\ & 42\quad 57 \\
& 28 & \ldots\ & 42\quad 14 \\
& Oct.\ 7 & \ldots\ & 42\quad 40 \\
& 10 & \ldots\ & 43\quad 15 \\
& 12 & \ldots\ & 42\quad 20 \\
& 20 & \ldots\ & 42\quad 5 \\
& 22 & \ldots\ & 42\quad 52 \\
\multicolumn{3}{l}{Mittel, Oct.\ 8} & $67^\circ\; 42'\; 48''$
\end{tabular}
\end{center}

Die ersten acht Beobachtungen sind auf ähnliche Art angestellt, wie die diesjäh-
rigen, indem an jedem Tage, ohne die Pole zwischen den Beobachtungen umzu-
kehren, zwei Nadeln (Nr.\ 1 und 2) angewandt wurden; die beiden letzten hinge-
gegen wurden auf die gewöhnliche Art gemacht, die zweite vom 20.\ Oct.\ mit Na-
del 4, die vom 22.\ mit Nadel 3. Die Zeit war am 27.\ Sept.\ und 10.\ Oct.\ Nach-
mittags zwischen 3 und 5 Uhr, bei allen übrigen Vormittags. Jede dieser 10 
Inclinationen beruhete auf 16 Einstellungen, und es wird ihnen aus diesem 
Grunde auch nur ein verhältnissmässig kleineres Gewicht zuzuerkennen sein, als 
den Inclinationen von 1842, die resp.\ auf 32, 64 und 40 Einstellungen beruheten.

\subsection*{33.}
\addcontentsline{toc}{subsection}{33.}

Sämmtliche bisher angeführte Inclinationen bedürfen noch einer kleinen 
gemeinschaftlichen Correction wegen des Einflusses, welchen an dem Beobach-
tungsplatze die Magnetstäbe der Magnetometer in der Sternwarte und im magne-
tischen Observatorium ausüben. Um die Resultate davon zu befreien, muss durch-
gehens $5{,}''15$ abgezogen werden (vergl.\ Resultate für 1840.\ II.\ Art.\ 6) [S.\,433 d.\,B.]
Die absolute Zuverlässigkeit der Inclinationsbestimmungen bleibt übrigens 
noch abhängig von der Richtigkeit der Voraussetzung, dass das Instrument selbst 
keine Theile enthält, die eine magnetische Wirkung auf die Nadel haben können. 
Ein Grund zu einer solchen Befürchtung ist bei dem von mir gebrauchten Instru-
mente nicht vorhanden; einige Beobachtungen, die ich nach der im 18.\ Art.\ er-
wähnten Art mit einer belasteten Nadel anstellte, haben immer nur Abweichun-
gen von ein Paar Minuten gezeigt, die sich aus den unvermeidlichen zufälligen 
Beobachtungsfehlern und den wirklichen Anomalien der Inclination selbst ganz 
ungezwungen erklären lassen. Auch die hinlänglich befriedigende Übereinstim-
mung der Werthe, welche im 11.\ Art.\ für die daselbst mit $\alpha$ bezeichnete Grösse 
gefunden sind, spricht gegen das Vorhandensein von solcher Störungen. Zur 
Erkennung ganz kleiner Einflüsse sind freilich solche Prüfungen nicht geeignet, 
und ich muss mir daher die weitere Prüfung durch mehr durchgreifende Mittel 
vorbehalten.

\subsection*{34.}
\addcontentsline{toc}{subsection}{34.}

Zum Schluss stelle ich noch meine Resultate mit einigen ältern Bestimmun-
gen zusammen.

\begin{center}
\begin{tabular}{r@{\ }l@{\ }r@{\ }l@{\ }l}
1805 & Dec.\ & $69^\circ$ & $29'$ & \ldots \\
1826 & Sept.\ & 68 & 29 & $26''$ \\
1837 & Juli & 1 & 47 & 20 \\
& & 7 & 53 & 30 \\
1841 & Oct.\ 8 & 67 & 42 & 43 \\
1842 & Juni 21 & 67 & 39 & 39 \\
\end{tabular}
\end{center}

Die beiden ersten Beobachtungen habe ich aus den Additions zu dem XIII.\ Bande 
der \textit{Voyage aux régions Equinoxiales} entlehnt (S.\,152); die erste ist mit einem In-
clinatorium von \textsc{Lenoir}, die zweite mit einem Instrument von \textsc{Gausey} angestellt; 
letztere beruhet auf den Beobachtungen mit zwei Nadeln, deren Resultate a.a.O.\ 
zu $68^\circ\; 30'\; 7''$ und $68^\circ\; 28'\; 15''$ angegeben werden, womit das ebendaselbst ange-
setzte Mittel nicht übereinstimmt; vermuthlich ist die Zahl für die zweite Nadel 
durch einen Druckfehler um $30''$ zu klein angesetzt. Der Beobachtungsplatz 1805 
ist mir nicht bekannt; 1826 war er im freien Felde einige hundert Schritte öst-
lich von der Sternwarte.

\textsc{Forbes}'s Beobachtungen sind in den \textit{Transactions of the Royal Society of Edin-
burgh} Vol.\,XV, Part.\,1, S.\,31 und 32 abgedruckt; sie wurden an einem \textsc{Robinson}-
schen Instrument von kleinern Dimensionen als das hiesige mit zwei Nadeln von 
6 engl.\ Zoll Länge im Garten der Sternwarte angestellt; die zweite Nadel hält der 
Beobachter selbst für die bessere.

Ich habe unter diese Beobachtungen die von \textsc{Mayer} im März 1814 ange-
stellten und in den \textit{Commentationes recent.\ Soc.\ Gotting.}\ T.\,III, S.\,36 u.\,37 ange-
führten nicht einreihen wollen, da dieselben gar kein Vertrauen verdienen. Wie 
sehr unvollkommen das von \textsc{Mayer} gebrauchte Instrument war, zeigt die von ihm 
selbst S.\,35 gegebene Probe, wo bei bleibender Stellung des Instruments zehn 
wiederholte Einstellungen Differenzen von mehr als einem Grade gaben. Seine 
Resultate für die Inclination selbst, von zwei verschiedenen Tagen, weichen um 
einen halben Grad von einander ab.

Eben so wenig verdient meine eigne Beobachtung vom 23.\ Juni 1832, die 
in der \textit{Intens.\ vis magneticae terrestris} art.\,27 angeführt ist, hier einen Platz, sowohl 
wegen der Unvollkommenheit des Instruments, als wegen des Locals in der Stern-
warte, wo nicht sehr entferntes Eisenwerk das Resultat bedeutend afficiren und 
zwar nachweislich eine Vergrösserung der Inclination hervorbringen musste.

Die angeführten Inclinationen lassen sich nun zwar sehr gut durch die An-
nahme einer jährlichen gleichförmigen Verminderung von 3 Minuten oder ge-
nauer $3'\; 2{,}''3$ vereinigen, wenn man bei \textsc{Forbes}'s Beobachtungen sich an das Re-
sultat der zweiten Nadel hält, und es bleiben nur Abweichungen übrig, die füg-
lich dem Conspiriren der Beobachtungsfehler und der Schwankungen der Incli-
nation zugeschrieben werden können. Da jedoch nach \textsc{Hansteen}'s Untersuchun-
gen über die Beobachtungen an andern europäischen Orten die jährliche Abnahme 
allmählig langsamer geworden ist, so wird man die angegebene Zahl nur wie ei-
nen mittlern etwa für 1829 gültigen Werth zu betrachten, und die Bestätigung 
und genauere Festsetzung der Ungleichförmigkeit erst von künftigen Beobachtun-
gen zu erwarten haben.

\newpage
%==============================================================================
% SECTION: Aufsätze über verschiedene Gegenstände der mathematischen Physik
% Page 505
%==============================================================================
\section*{Aufsätze über verschiedene Gegenstände der mathematischen Physik}
\addcontentsline{toc}{section}{Aufsätze über verschiedene Gegenstände der mathematischen Physik}

\begin{center}
\Large\bfseries
Aufsätze\\[0.5em]
über verschiedene Gegenstände\\[0.5em]
der mathematischen Physik.
\end{center}

\newpage
%==============================================================================
% SECTION: Fundamentalgleichungen für die Bewegung schwerer Körper auf der rotirenden Erde
% Pages 507-515
%==============================================================================
\section*{Fundamentalgleichungen für die Bewegung schwerer Körper auf der rotirenden Erde}
\addcontentsline{toc}{section}{Fundamentalgleichungen für die Bewegung schwerer Körper auf der rotirenden Erde}

\textsc{Benzenberg.}\ \textit{Versuche über das Gesetz des Falls.}\ 1804.\ Seite 349 bis 353 und 363 bis 371.

Brief von Gauss an \textsc{Benzenberg}.

Braunschweig 1803, Februar 2.

\medskip

\noindent
--- In der Theorie unsres Freundes \textsc{Olbers} ist eine Voraussetzung, die 
mir nicht zulässig scheint. Nemlich: dass der Körper während des Falls in einer 
Ebene bleibe. Allein dies darf man, meiner Meinung nach, nicht voraussetzen, 
wenn man den Widerstand der Luft in Betracht zieht, den man hier nothwendig 
in Betracht ziehen muss, weil die geschlossene Abweichung nach Süden lediglich 
darauf beruht. Eine leichte Betrachtung zeigt nemlich folgendes: die Ebene ($A$), 
in welcher der Körper sich ursprünglich zu bewegen anfängt, geht durch den Mit-
telpunkt der Erde (oder allgemeiner, der Attraction), und steht auf derjenigen 
Ebene ($B$) senkrecht, in der der Meridian des Beobachtungsorts beim Anfang 
des Falls war. Allein man sieht leicht, dass die Lufttheile an allen Stellen der 
Ebene $A$ schief durchgehen, bloss die gerade Linie ausgenommen, wo $A$ von $B$ geschnitten wird. Die Luft wirkt daher dem Körper nicht in dieser Ebene $A$ 
entgegen, sondern treibt ihn daraus weg nach Norden, und es schien mir, dass 
der Effect davon gerade so gross sein würde, dass er die aus der Verspätung des 
Falls geschlossene Abweichung nach Süden aufhöbe.

Nachdem ich durch Ihren letzten Brief veranlasst war, aufs Neue an diese 
Materie zu denken, betrachtete ich in einer müssigen halben Stunde die Sache 
auf eine ganz verschiedene Art, und entwickelte die analytischen Gleichungen, 
die die relative Bewegung des Körpers gegen die bewegte Erdoberfläche in sich 
fassen, aus den ersten Fundamentalsätzen der Dynamik, und hier fand ich zu 
meiner Verwunderung

\begin{enumerate}
\item die Abweichung nach Süden wiederum 0 oder ganz unvermerklich;
\item die Abweichung nach Osten nur $\tfrac{7}{8}$ von dem, was Dr.\ \textsc{Olbers} gefunden 
hat. Nemlich in Dr.\ \textsc{Olbers}' Zeichen, wenn man den Widerstand der Luft ver-
nachlässigt,
\[
y = \frac{2n\cos\varphi\cdot a}{86164}
\]
oder wenn man ihn mit in Betrachtung zieht, nach einer hier zureichenden Nä-
herung
\[
y = \frac{2n\cos\varphi\cdot a}{7{,}896164\,(\delta a - \delta' a)}
\]
\end{enumerate}

wo $a$ die Höhe ist, durch die der Körper in der Zeit $t$ im leeren Raume fallen 
würde, also $= \tfrac{1}{2}gt$ [wo ferner $\delta$ die Polhöhe des Beobachtungsortes und $a'$ 
die wirkliche Fallhöhe bezeichnet].

Hienach finde ich für Ihre Versuche, indem ich die Pendellänge für Ham-
burg $= 440{,}75$ Linien (woraus $g$ fast eben so kommt, wie Dr.\ \textsc{Olbers} 
\ldots
giebt), also 
\[
g = 15{,}62925\text{ Fuss (Hamburg)}
\]
und die wirkliche Fallhöhe bei dem Luftwiderstande
\[
\delta a - \delta' a = 235{,}51\text{ Fuss,}
\]
die Abweichung nach Osten 
\[
y = 0{,}10805\text{ Fuss} = 12{,}966\text{ Linien,}
\]
während Sie gefunden haben 14,04 Linien. Die Abweichung nach Süden ist 
ganz unmerklich.

\newpage

\textsc{Gau\ss.}

\begin{center}
\Large\bfseries
Fundamentalgleichungen\\[0.5em]
für die Bewegung schwerer Körper\\[0.5em]
auf der rotirenden Erde.
\end{center}

\medskip

Ich bezeichne die Erdaxe als die $z$, ihren entgegengesetzten Richtungssinn 
als $z'$, die Ebene des Meridians des Beobachtungsorts (oder vielmehr die paral-
ele durch den Anfangspunkt der relativen Coordinaten) als die $zx$, und die 
entgegengesetzte Richtung von $x$ als $x'$; die Richtung der $y$ ist so zu wählen, 
dass $xyz$ rechtwinklig sind. In Beziehung auf diese Axen sind die Coordinaten 
eines Lufttheilchens
\[
x,\quad y,\quad z
\]
und eines Erdelementes
\[
a + X,\quad b + Y,\quad c + Z,
\]
wo $a$, $b$, $c$ die Coordinaten des Mittelpunkts der Attraction sind. Ich setze
\begin{align*}
X &= (z-a)\sin p\cos\vartheta + y\sin\vartheta + (z-c)\cos p\cos\vartheta,\\
Y &= -(z-a)\sin p\sin\vartheta + y\cos\vartheta - (z-c)\cos p\sin\vartheta,\\
Z &= -(z-a)\cos p \qquad\qquad + (z-c)\sin p.
\end{align*}

Die Coordinaten $X$, $Y$, $Z$ lassen sich einerseits als Functionen von $t$ al-
lein, anderseits aber auch als Functionen der vier veränderlichen Grössen $\vartheta$, $x$, $y$, $z$ 
betrachten, und haben also in letzterer Hinsicht vier partielle Differentiale. Es 
ist demnach
\begin{align*}
dX &= \tfrac{\partial X}{\partial\vartheta}d\vartheta + \tfrac{\partial X}{\partial x}dx + \tfrac{\partial X}{\partial y}dy + \tfrac{\partial X}{\partial z}dz,\\
dY &= \ldots,\\
dZ &= \ldots
\end{align*}

Die Geschwindigkeit des Körpers zerlegt sich, wie seine Bewegung, in drei 
partielle auf die Ebnen der $X$, $Y$, $Z$ senkrechte Geschwindigkeiten, die mithin
\[
\frac{dX}{dt},\quad \frac{dY}{dt},\quad \frac{dZ}{dt}
\]
sind. Die Geschwindigkeiten des Luftelements hingegen, in welchem er sich 
jedesmal befindet, in Beziehung auf dieselben Ebnen sind offenbar
\[
\frac{dX}{dt} - \frac{\partial X}{\partial x}\frac{dx}{dt} - \frac{\partial X}{\partial y}\frac{dy}{dt} - \frac{\partial X}{\partial z}\frac{dz}{dt} = -\sin p\cos\vartheta\cdot\tfrac{dx}{dt} - \sin\vartheta\cdot\tfrac{dy}{dt} + \cos p\cos\vartheta\cdot\tfrac{dz}{dt}.
\]

Folglich die relativen Geschwindigkeiten des Körpers nach diesen drei Richtungen
\begin{align*}
\xi &= \cos p\cos\vartheta\cdot\tfrac{dx}{dt} - \sin\vartheta\cdot\tfrac{dy}{dt} + \sin p\cos\vartheta\cdot\tfrac{dz}{dt},\\
\eta &= \cos p\sin\vartheta\cdot\tfrac{dx}{dt} + \cos\vartheta\cdot\tfrac{dy}{dt} + \sin p\sin\vartheta\cdot\tfrac{dz}{dt},\\
\zeta &= -\sin p\cdot\tfrac{dx}{dt} \qquad\qquad + \cos p\cdot\tfrac{dz}{dt}.
\end{align*}

Die totale relative Geschwindigkeit ist folglich $= \sqrt{(\xi\xi + \eta\eta + \zeta\zeta)}$, welche, 
wie die Entwickelung aus obigen Werthen leicht zeigt, $= \sqrt{(\dot{x}^2 + \dot{y}^2 + \dot{z}^2)}$ 
wird. Der Widerstand der Luft ist dem Quadrate davon proportional; wir 
setzen ihn daher $= Muu$, und zerlegen ihn nach obigen drei Richtungen in 
$M\mu\xi$, $M\mu\eta$, $M\mu\zeta$.

Wir sehen hier die Erde als ein Revolutions-Sphäroid an; die Richtung der 
Schwere geht daher durch die Erdaxe. Der Punkt, wo sie diese schneidet, liege 
um $g$ über $C$, oder es sei für denselben $Z = g$.

Setzt man nun ferner die Stärke der Gravitation $= \gamma$ und
\[
\gamma(z-a) = P,\quad \gamma y = Q,\quad \gamma(z-c) = R,
\]
so ist nach den Grundsätzen der Dynamik
\begin{align}
0 &= \frac{d^2X}{dt^2} + \tfrac{\partial X}{\partial x}P + \tfrac{\partial X}{\partial y}Q + \tfrac{\partial X}{\partial z}R + M\mu\xi,\\
0 &= \frac{d^2Y}{dt^2} + \tfrac{\partial Y}{\partial x}P + \tfrac{\partial Y}{\partial y}Q + \tfrac{\partial Y}{\partial z}R + M\mu\eta,\\
0 &= \frac{d^2Z}{dt^2} + \tfrac{\partial Z}{\partial x}P + \tfrac{\partial Z}{\partial y}Q + \tfrac{\partial Z}{\partial z}R + M\mu\zeta.
\end{align}

Aus obigen Werthen von $X$, $Y$, $Z$ findet man, wenn man für 
$\tfrac{d\vartheta}{dt}$, welches beständig ist, $n$ schreibt, folgende Gleichungen:
\begin{align}
\frac{d^2X}{dt^2} &= \ddot{x}\sin p\cos\vartheta + \ddot{y}\sin\vartheta + \ddot{z}\cos p\cos\vartheta \notag\\
&\quad - 2n\dot{x}\sin p\sin\vartheta - 2n\dot{y}\cos\vartheta - 2n\dot{z}\cos p\sin\vartheta - n^2X,\\
\frac{d^2Y}{dt^2} &= -\ddot{x}\sin p\sin\vartheta + \ddot{y}\cos\vartheta - \ddot{z}\cos p\sin\vartheta \notag\\
&\quad - 2n\dot{x}\sin p\cos\vartheta + 2n\dot{y}\sin\vartheta + 2n\dot{z}\cos p\cos\vartheta - n^2Y,\\
\frac{d^2Z}{dt^2} &= -\ddot{x}\cos p \qquad\qquad + \ddot{z}\sin p \notag\\
&\quad + 2n\dot{x}\cos p - 2n\dot{z}\sin p.
\end{align}

Multiplicirt man die drei Gleichungen (3) resp.\ mit $\sin p\cos\vartheta$, $\sin p\sin\vartheta$, $-\cos p$ 
und addirt die Producte; multiplicirt man zweitens eben diese Gleichungen mit 
$-\sin\vartheta$, $\cos\vartheta$, $0$; und drittens mit $\cos p\cos\vartheta$, $\cos p\sin\vartheta$, $\sin p$, und addirt bei-
demale die Producte: so erhält man, nachdem man statt $\tfrac{d^2X}{dt^2}$, $\tfrac{d^2Y}{dt^2}$, $\tfrac{d^2Z}{dt^2}$ ihre 
Werthe aus (4), statt $X$, $Y$, $Z$ die aus (2), und statt $\xi$, $\eta$, $\zeta$ die ihrigen substi-
tuirt hat, folgende drei neue:
\begin{align}
0 &= \ddot{x} - 2n\sin p\,\dot{y} + (x-a)(n^2 - \gamma\mu) + \cos p\,\gamma\mu(Z-z) + M\mu\xi,\\
0 &= \ddot{y} + 2n\sin p\,\dot{x} - 2n\cos p\,\dot{z} + y(n^2 - \gamma\mu) + M\mu\eta,\\
0 &= \ddot{z} + 2n\cos p\,\dot{y} + (z+c)(n^2 - \gamma\mu) - \sin p\,\gamma\mu(Z-z) + M\mu\zeta.
\end{align}

Ist also der Körper gegen die Erde in Ruhe, oder $\dot{x} = \dot{y} = \dot{z} = 0$, so 
scheint er senkrecht auf die Ebnen der $x$, $y$, $z$ von den Kräften
\begin{align*}
&(x-a)(n^2 - \gamma\mu) + \cos p\,\gamma\mu(Z-z),\\
&y(n^2 - \gamma\mu),\\
&(z+c)(n^2 - \gamma\mu) - \sin p\,\gamma\mu(Z-z)
\end{align*}
sollicitirt zu werden. Ein schon in Bewegung begriffener Körper hingegen wird 
anders afficirt. Denn ausser dem Widerstande der Luft, der den Körper nach 
diesen Richtungen wie Kräfte, deren Maass $M\mu\xi$, $M\mu\eta$, $M\mu\zeta$ ist, treibt 
und folglich auf der rotirenden Erde völlig eben so wirkt, als er in der ruhen-
den wirken würde, kommen nach jenen Richtungen noch die drei Kräfte
\[
-2n\sin p\,\dot{y},\quad +2n\sin p\,\dot{x} - 2n\cos p\,\dot{z},\quad +2n\cos p\,\dot{y}
\]
hinzu, und diese sind es allein, wodurch die Rotation der Erde an fallenden Kör-
pern sichtbar wird. Die bisherigen Schlüsse und Folgerungen sind streng und 
allgemein richtig.

Bei Versuchen, die in dieser Hinsicht angestellt werden, geschieht allemal 
die Bewegung des Körpers in einem so kleinen Raume, dass man die Stärke der 
auf ruhende Körper wirkenden scheinbaren Schwere innerhalb desselben, als un-
veränderlich $= g$, und ihre Richtung als immer parallel, also senkrecht auf die 
Ebne der $z$ annehmen kann. Es wird also ohne Bedenken erlaubt sein, statt 
der obigen drei Grössen
\begin{align*}
&(x-a)(n^2 - \gamma\mu) + \cos p\,\gamma\mu(Z-z),\\
&y(n^2 - \gamma\mu),\\
&(z+c)(n^2 - \gamma\mu) - \sin p\,\gamma\mu(Z-z)
\end{align*}
respective $0$, $0$, $g$ zu substituiren. Dadurch werden die drei Fundamentalglei-
chungen
\begin{align}
0 &= \ddot{x} - 2n\sin p\,\dot{y} + M\mu\xi,\\
0 &= \ddot{y} + 2n\sin p\,\dot{x} - 2n\cos p\,\dot{z} + M\mu\eta,\\
0 &= \ddot{z} + 2n\cos p\,\dot{y} - g + M\mu\zeta.
\end{align}

Die Integration dieser Gleichungen ist leicht, wenn man den Widerstand 
der Luft vernachlässigt, oder $M = 0$ setzt. Man findet nemlich
\begin{align*}
x &= A - D\cos p\cdot t + E\sin p\cos(2nt + \varepsilon) + \tfrac{1}{2}\sin p\cos p\cdot gt^2,\\
y &= B - \tfrac{1}{2}E\sin(2nt + \varepsilon) + \tfrac{1}{2}\cos p\cdot gt,\\
z &= C + D\sin p\cdot t - E\cos p\cos(2nt + \varepsilon) - \tfrac{1}{2}\sin p^2\cdot gt^2.
\end{align*}

Auch ist es leicht, folgende Werthe der arbiträren Grössen zu entwickeln, 
wenn man voraussetzt, dass der Körper anfänglich gar keine scheinbare Ge-
schwindigkeit hat:
\[
A = -c\sin p,\quad B = b,\quad C = -c\cos p,\quad D = b\cos p,\quad E = c,\quad F = \tfrac{g}{4n^2}.
\]

Also
\begin{align*}
x &= c\cos p\sin p\,(ntt - \tfrac{1}{2} + \tfrac{1}{2}\cos 2nt),\\
y &= \tfrac{1}{2}c\cos p\,(1 - \sin 2nt),\\
z &= -\tfrac{1}{2}gtt + c\cos p^2\,(ntt - \tfrac{1}{2} + \tfrac{1}{2}\cos 2nt).
\end{align*}

Diese Integration ist freilich nicht allgemein zulässig, da obige Voraus-
setzung nur in so fern erlaubt ist, als der Körper sich von seinem anfänglichen 
scheinbaren Orte nicht weit entfernt. Für diesen Fall aber können wir die tri-
gonometrischen Functionen in Reihen auflösen, und so wird
\begin{align*}
x &= \tfrac{1}{2}c\cos p\sin p\cdot gnnt^4 \ldots\\
y &= \tfrac{1}{4}t\cos p\cdot gnt^3 \ldots\\
z &= -\tfrac{1}{2}gtt + \tfrac{1}{4}t\cos p^2\cdot gnnt^4 \ldots
\end{align*}

Da die Zeit des Falls nur wenige Secunden, also $nt$ höchstens einige Raum-
minuten beträgt, und (weil Radius $= 1$) $1' = \tfrac{1}{3438}$, so wird $x$ und der zweite 
Theil von $z$ ganz unmerklich, also $y = -\tfrac{1}{3}z\cos p\cdot nt$. Bei Dr.\ \textsc{Benzenberg}'s 
Versuche im Michaelisthurm war $z = -235$ Fuss, $\varphi = 53^\circ\,33'$, $t = 4''$ Son-
enzeit, also $nt = \tfrac{4}{86164}$ Raumminuten. Hieraus wird $y = 3{,}91$ Linien.

Wenn man bei der Integration obiger Gleichungen den Widerstand der 
Luft mit in Betrachtung ziehen will, so wird man sich mit Näherungen begnü-
gen müssen; die Entwickelung der Werthe von $x$, $y$, $z$ in Reihen nach den Po-
tenzen von $n$ und $M$ ist alsdann sehr leicht. Das höchste Glied von $x$ wird 
wie vorhin $= \tfrac{1}{2}c\cos p\sin p\cdot gnnt^4$, und ist also von gar keiner Bedeutung; für 
$y$ und $z$ findet man mit Vernachlässigung der Quadrate und höhern Potenzen 
von $n$ und $M$ folgende Werthe:
\begin{align*}
y &= \tfrac{1}{4}t\cos p\cdot gnt^3 - \tfrac{1}{720}\cos p\cdot Mggnt^6,\\
z &= -\tfrac{1}{2}gtt + \tfrac{1}{24}Mggt^4.
\end{align*}

Setzen wir also $-z$, den wirklichen Fall, $= f$; $\tfrac{1}{2}gtt$ oder den Fall im 
luftleeren Raume $= f + \delta$, so ist
\[
y = \tfrac{1}{3}\cos p\cdot nt(f + \delta) - \cos p\cdot nt\delta = \tfrac{1}{3}\cos p\cdot nt(f - 2\delta).
\]

Für die Versuche in St.\ Michael, wo $f + \delta = 241{,}47$ Fuss war, erhal-
ten wir daher die Abweichung nach Osten $y = 3{,}86$ Linien.

\newpage
%==============================================================================
% SECTION: Über die achromatischen Doppelobjective
% Pages 516-522
%==============================================================================
\section*{Über die achromatischen Doppelobjective besonders in Rücksicht der vollkommenern Aufhebung der Farbenzerstreuung}
\addcontentsline{toc}{section}{Über die achromatischen Doppelobjective}

\textit{Zeitschrift für Astronomie und verwandte Wissenschaften herausgegeben von 
B.\ von \textsc{Lindenau} und \textsc{Bohnenberger}. Bd.\ IV. N.\ XXX. 1817. December. Seite 345 bis 351.}

Der schöne Aufsatz des Hrn.\ Prof.\ \textsc{Bohnenberger} über die achromatischen 
Objective im ersten Bande dieser Zeitschrift hat das Verdienst, einen für diese 
Theorie wichtigen Umstand zuerst zur Sprache gebracht zu haben. Ich bin da-
durch veranlasst, einige frühere Untersuchungen wieder vorzunehmen und wei-
ter zu entwickeln, deren Resultate ich hier mittheilen werde.

Man begnügte sich bisher bei den Doppelobjectiven, die Farbenzerstreuung 
für die der Axe unendlich nahen Strahlen, und die Abweichung wegen der Ku-
gelgestalt für die Strahlen von mittlerer Brechbarkeit zu heben, wobei also für 
die Randstrahlen noch eine kleine Farbenzerstreuung zurückbleiben kann. Bei 
dieser Einrichtung ist die Berechnung des achromatischen Objectivs eine unbe-
stimmte Aufgabe, d.i., zu jeder Kronglaslinse von positiver Brennweite, wie 
auch immer das Verhältniss der Halbmesser der Flächen sein mag, lässt sich 
eine Flintglaslinse berechnen, die mit jener vereinigt ein in obiger Bedeutung 
achromatisches Objectiv gibt. So viel ich weiss, haben bisher alle Optiker beide 
Flächen der Kronglaslinse convex angenommen: allein für das Verhältniss der 
beiden Halbmesser haben die Theoretiker sehr verschiedene Werthe in Vorschlag 
gebracht, je nachdem sie von diesem oder jenem Princip ausgingen. Will man 
mit \textsc{Euler} die Abweichung wegen der Gestalt bei der Kronglaslinse zu einem 
Kleinsten machen, so müssen die Halbmesser ungefähr in dem Verhältniss von 
1 zu 7 stehen; sie müssen einander gleich sein, wenn man, wie \textsc{Kügel} in der 
analytischen Dioptrik, die möglich kleinsten Krümmungen zu haben wünscht; 
sollen die Brechungen selbst die möglich kleinsten werden, wie derselbe Schrift-
steller in einer spätern Abhandlung sich vorsetzt, so müssen diese Brechungen 
einander gleich sein, und die Halbmesser nahe in dem Verhältniss von 1 zu 3 
stehen. Es scheint nicht, dass alle diese verschiedenen Vorschläge hinlänglich 
motivirt sind. \textsc{Kügel}'s Augenmerk ist besonders die Abweichung wegen der Ku-
gelgestalt gewesen, welche für ale Strahlen in mathematischer Schärfe zu heben 
bekanntlich unmöglich ist: bei \textsc{Euler}'s Behandlung dieser Rechnungen ist diese 
Abweichung eigentlich nur für die der Axe nächsten Strahlen gehoben, und es 
bleibt eine sehr nahe dem Biquadrat des Abstandes von der Axe proportionale, 
also für die Randstrahlen am meisten merkliche Abweichung zurück; oder wenn 
man mit \textsc{Kügel} die Rechnung so führt, dass die Abweichung für die Randstrah-
len verschwindet, so kommt sie wieder bei den Zwischenstrahlen zum Vorschein, 
am merklichsten bei denen, deren Entfernung nahe $\tfrac{1}{2}$ von dem Halbmesser der 
Öffnung ist. Die unvermeidlich übrig bleibende Abweichung wegen der Gestalt 
so unschädlich wie möglich zu machen, war \textsc{Kügel}'s Absicht bei der Wahl des 
Verhältnisses der beiden ersten Halbmesser; es erhellt jedoch nicht klar genug 
weder, dass wirklich dieser Zweck bei dem gewählten Verhältniss am allerbesten 
erreicht werde, noch, dass dieser Zweck wichtig genug sei, um ihn vorzugsweise 
allein zur Grundlage der Bestimmung dieses Verhältnisses zu machen. Finden 
nemlich noch andere Umstände Statt, die mit jener Abweichung in nahem Zu-
sammen-
\ldots

Das Resultat meiner Rechnung ist nun folgendes:

Wenn die Halbmesser der Reihe nach zu
\[
+3415{,}287;\quad -10133{,}007;\quad -4207{,}421;\quad -2307{,}320
\]
angenommen werden, so vereinigen sich die rothen und violetten Strahlen, so-
wohl die, welche unendlich nahe bei der Axe, als die, welche in der Entfernung 
$1083{,}687$ auffallen, alle in einem Punkt der Axe, dessen Entfernung von der 
letzten Fläche $= 28293{,}3$ wird. Wird jene Entfernung von der Axe, bei wel-
der der Einfallswinkel $18^\circ\,30'$ ist, als Halbmesser der Öffnung angenommen, 
so ist der Durchmesser der Öffnung sehr nahe $\tfrac{1}{4}$ der Brennweite. Um beurthei-
len zu können, wie gross die noch übrig bleibende Abweichung wegen der Ge-
stalt für die Strahlen zwischen dem Rande und der Axe wird, habe ich die Ver-
einigungsweiten für den Einfallswinkel $13^\circ$ berechnet und gefunden
\begin{align*}
&28289{,}3\text{ für die rothen}\\
&28290{,}0\text{ für die violetten Strahlen.}
\end{align*}

Ich kann nicht umhin, hier noch eine Erinnerung über eine Äusserung des 
Hrn.\ Prof.\ \textsc{Bohnenberger} in dem erwähnten Aufsatze beizufügen. Ich halte nem-
lich dafür, dass es am vortheilhaftesten ist, die Abweichung wegen der Gestalt 
genau für die Randstrahlen zu heben. Hr.\ Prof.\ \textsc{Bohnenberger} hat S.\,279 dieses 
Verfahren wie mir deucht mit Unrecht getadelt. Man könnte, sagt er, wenn die 
Abweichung für die Randstrahlen genau gehoben sei, die Öffnung ohne Schaden 
der Deutlichkeit bis dahin vergrössern, wo die Abweichung wieder der grössten 
Abweichung der Zwischenstrahlen gleich werde, und es sei daher am vortheil-
haftesten, die Abweichung nicht für die Randstrahlen, sondern für Strahlen zwi-
schen dem Rande und der Axe zu heben. Dies würde allerdings wahr sein, 
wenn die übrigbleibende Abweichung jenseits und diesseits der Entfernung, für 
welche sie gehoben ist, einerlei Zeichen hätte, was aber nicht der Fall ist. Man 
könnte zwar hiegegen mit einigem Schein einwenden, dass es bei der Längenab-
weichung auf das Zeichen gar nicht ankomme, und dass positive und negative 
Abweichungen eine und dieselbe Undeutlichkeit im Auge hervorbringen. Allein 
hiebei nähme man offenbar stillschweigend an, dass das Ocular immer genau für 
das deutliche Sehen desjenigen Bildes gestellt sei, welches durch die der Axe 
nächsten Strahlen hervorgebracht wird, und dies kann doch nicht eingeräumt 
werden. Man mag dies Bild immerhin das Hauptbild nennen: es fällt mit dem 
von den Randstrahlen hervorgebrachten Bilde zusammen, wenn die Abweichung 
für diese gehoben ist, und alle übrigen Bilder werden dann (wenigstens allgemein 
zu reden) jenseits oder diesseits des Hauptbildes liegen. Da man nun das Ocu-
lar immer so stellt, dass die Undeutlichkeit so klein wie möglich wird, so sieht 
man gerade das Hauptbild am wenigsten deutlich, und jede Vergrösserung der 
Öffnung vergrössert auch die Undeutlichkeit. Eine ausführlichere Erörterung 
dieses Umstandes würde mich hier zu weit abführen.

\subsection*{[Brief von Gauss an Brandes]}
\addcontentsline{toc}{subsection}{Brief von Gauss an Brandes}

\textsc{Gehler}'s \textit{Physikalisches Wörterbuch.} 1831. Artikel: 
Linsenglas, Berechnungen über achromatische und aplanatische Linsengläser aus zwei Glaslinsen.

Brief von Gauss an \textsc{Brandes}.

Auf Veranlassung Ihres Briefes habe ich eine freie Stunde auf den in jenem 
Aufsatze\footnote{Über die achromatischen Doppelobjective besonders in Rücksicht der vollkommenern Aufhebung der Farbenzerstreuung.} am Ende kurz erwähnten Umstand gewandt. Der eigentliche Sinn 
der dortigen Bemerkung scheint nicht von allen ganz richtig aufgefasst zu sein, 
aber auch meine Angabe bedarf einer kleinen Modification. Ich finde nemlich 
jetzt durch eine tiefer eindringende Untersuchung, dass die Undeutlichkeit, die 
in dem Ausdrucke für die Längen-Abweichung von der vierten Potenz des Abstan-
des der auffallenden Strahlen von der Axe abhängt, den möglich kleinsten To-
tal-Einfluss hat, wenn man das Objectiv so construirt, dass diejenigen Strahlen, 
die unendlich nahe bei der Axe einfallen, und diejenigen, die in einer Entfernung 
$= R\sqrt{\tfrac{3}{5}}$ auffallen würden, (wo $R =$ Radius des Objectivs ist) in einem 
Punkte $A$ sich vereinigen, wobei das Ocular dann so steht, dass man denjeni-
gen Punkt der Axe, wo die Strahlen, die in der Entfernung $= CR$ und 
$= C + \delta R$ von der Axe aufgefallen sind, sich alle vereinigen, deutlich sieht.

Denken Sie Sich nemlich durch diesen Punkt eine auf die Axe senkrechte Ebene, 
so ist das Bild desto undeutlicher, je grösser der Kreis um $A$ ist, den die von 
einem Punkte des Objects auf das Objectivglas gefallenen Strahlen füllen, doch 
so, dass die Intensität der Strahlen an jeder Stelle dieses Kreises mit berücksich-
tigt werden muss. Hiebei ist nun einige Willkürlichkeit; ich halte für das 
zweckmässigste, hier nach denselben Principien zu verfahren, die der Methode 
der kleinsten Quadrate zum Grunde liegen. Ist nemlich $ds$ ein Element dieses 
Kreises, $\rho$ die Entfernung des Elements von $A$, und $i$ die Intensität der Strah-
len daselbst, so nehme ich an, dass $\int i\rho\rho\,ds$ als das Maass der Total-Undeutlich-
keit zu betrachten sei, und mache dies zu einem Minimum. Ich finde dabei 
folgende Resultate:

\begin{enumerate}
\item Construirte man das Objectiv so, dass dasjenige Glied der Längen-Abweichung, welches von dem Quadrate der Entfernung von der Axe abhängt, $= 0$ wird, und setzte das Ocular so, dass $A$ dahin fällt, wo die der Axe unendlich nahen Strahlen diese schneiden, so sei der Werth dieses Integrals $= E$.
\item Stellte man aber bei derselben Einrichtung das Ocular so, dass das Integral so klein wird, wie es bei dieser Einrichtung werden kann (wobei $A$ der Vereinigungspunkt der in der Entfernung $= R\sqrt{\tfrac{1}{2}}$ auffallenden Strahlen sein wird), so ist das Integral $= \tfrac{5}{8}E$.
\item Dagegen ist bei der obigen Einrichtung und der vortheilhaftesten Stellung des Oculars das Integral $= \tfrac{4}{9}E$, als absolutes Minimum.
\end{enumerate}

Obiges Resultat, dass nemlich mit dem Vereinigungspunkte der der Axe unendlich nahen Strahlen ein bloss fingirtes Bild (von Strahlen aus grösserer Distanz von der Axe als der Halbmesser des Objectivs) vereinigt werden soll, ist anfangs sehr überraschend und paradox scheinend; aber bei näherer Be-
trachtung sieht man den eigentlichen Grund leicht ein. Jenes erste sogenannte 
Hauptbild (von Strahlen sehr nahe bei der Axe) ist nemlich dabei gleichsam das 
Unwichtigste wegen seiner geringen Intensität, viel wichtiger ist, dass die Strah-
len von den der Peripherie näheren Ringen des Objectivs unter sich besser zusam-
men gehalten werden, was bei jener Einrichtung am besten erreicht wird. Es 
thut mir leid, dass die Grenzen eines Briefes jetzt grössere Ausführlichkeit nicht 
gestatten; der scharfe Calcül lässt sich nichts abstreifen und bei einem vagen Rai-
sonnement übersieht man leicht einen wesentlichen Umstand; allein für den Ken-
ner werden diese Winke schon zureichen.

Allgemein finde ich, dass immer bei der vortheilhaftesten Stellung des Ocu-
lars jenes Integral $= \tfrac{1}{9}E(1 - \tfrac{6}{5}pp + p^4)$ wird, wenn das Objectiv so construirt 
ist, dass Strahlen aus der Entfernung $pR$ von der Axe sich mit dem (oben so-
genannten) Hauptbilde in einem Punkte vereinigen. Dies ist ein Minimum 
für $p = \sqrt{\tfrac{3}{5}}$ und ist dann $= \tfrac{4}{9}E$; für $p = 1$ wäre es nur $= \tfrac{5}{9}E$ und für 
$p = \infty$ undendlich klein, $= E$. Nicht allein hat also hienach \textsc{Bohnenberger} 
Unrecht, sondern auch ich habe damals Unrecht gehabt, aber insofern, als ich 
noch nicht weit genug von \textsc{Bohnenberger} abgewichen bin. Ich hatte damals bloss 
die ganze Grösse des undeutlichen Bildes berücksichtigt, ohne auf die ungleiche 
Intensität der einzelnen Theile Rücksicht zu nehmen.

\newpage
%==============================================================================
% SECTION: Berichtigung der Schneiden einer Waage
% Pages 523-526
%==============================================================================
\section*{[Berichtigung der Schneiden einer Waage]}
\addcontentsline{toc}{section}{Berichtigung der Schneiden einer Waage}

\textit{Göttingische gelehrte Anzeigen. Stück 41. Seite 401 bis 405. 1837 März 13.}

In der Sitzung der Königl.\ Gesellschaft der Wissenschaften vom 28.\ Ja-
nuar nahm der Hofr.\ Gauss von der Vorlesung des Hrn.\ Prof.\ \textsc{Weber}, über wel-
che im 22.\ Stücke dieser Blätter Bericht abgestattet ist, Veranlassung, einen Vor-
trag über einen nahe verwandten Gegenstand zu halten, von welchem wir den 
Hauptinhalt hier zur Anzeige bringen.

Er betrifft eine neue Berichtigungsmethode zur Erfüllung einer wesentli-
chen Bedingung bei den feineren Hebelwaagen, deren Wichtigkeit bisher nicht 
genug gewürdigt zu sein scheint. Solche Waage haben drei prismatische Schnei-
den; die eine nach unten gekehrte, in der Mitte des Waagebalkens, ruht auf 
einem harten horizontalen Lager von Stein oder Stahl, und dient als Drehungs-
axe bei dem Spiel des Waagebalkens; die beiden andern an den Enden des Waa-
gebalkens sind aufwärts gerichtet, und auf jeder derselben schwebt das Trage-
stück, woran die Waageschale hängt. Die Tragestücke selbst sind von gehärte-
tem Stahl, und ihre unteren, auf den Schneiden aufliegenden Flächen vollkom-
men plan und hochpolirt.

Eine wesentliche Bedingung ist nun, dass diese beiden äussern Schneiden 
mit der mittleren parallel sein sollen. In der That, da vor jedem Umtausch der 
Gewichte in einer Schale die Waage erst gehemmt und dabei das Tragestück von 
der Schneide abgehoben wird, so ist nie darauf zu rechnen, dass sich nach Auf-
hebung der Hemmung das Tragestück genau wieder eben so auf die Schneide 
legt, wie zuvor: dies ist zwar unschädlich, wenn die betreffende Schneide mit 
der mittleren parallel ist, verursacht aber ein verändertes Moment, wenn eine 
Divergenz der Schneiden statt findet. Eine unvollkommene Berichtigung in die-
ser Beziehung ist eine Hauptursache, warum bei oft wiederholten Wägungen zu-
weilen bedeutend grössere Abweichungen in den Resultaten sich zeigen, als man 
sonst von der vortrefflichen Arbeit und der Empfindlichkeit einer Waage erwar-
ten sollte.

Die Mittel, deren sich die Künstler zur Berichtigung des Parallelismus der 
Schneiden bisher gewöhnlich bedient haben, sind nicht geeignet, alle zu wün-
schende Schärfe zu geben; auch ist es, bei feinen Waagen wie bei astronomischen 
Instrumenten, nicht der Verfertiger, von dem man die feinste Berichtigung zu 
fordern hat, sondern diese kommt dem zu, der die Waage gebraucht.

Das Verfahren, dessen sich der Hofr.\ Gauss zu dieser Berichtigung mit 
dem besten Erfolge bedient hat, beruht auf folgenden Gründen.

Bei den Schwingungen des Waagebalkens verändert die zu prüfende äussere 
Schneide zwar ihre Lage im Raume; diese verschiedenen Lagen sind aber alle 
unter einander parallel, wenn diese Schneide mit der (ruhenden) mittleren paral-
llel ist. Anders verhält es sich dagegen, wenn die äussere Schneide der mittleren 
nicht parallel ist. Nehmen wir, um die Vorstellung zu fixiren, an, dass die 
äussere Schneide zwar mit der mittleren in Einer Ebene liege, dass aber die Rich-
tungen der beiden Schneiden abwärts vom Beobachter divergiren. In diesem 
Falle wird bei dem Spiele des Waagebalkens die äussere Schneide sich auf einer 
Kegelfläche bewegen; ihr abwärts gekehrtes Ende wird, relativ gegen das nähere 
Ende, steigen oder sinken, so wie der Hebelarm, an welchem diese Schneide 
sich befindet, steigt oder sinkt. Dasselbe wird von dem die Schneide stets be-
rührenden Tragestücke gelten.

Welcher von beiden Fällen nun statt finde, lässt sich erkennen, wenn auf 
dem Tragestücke ein Planspiegel befestigt ist. Am vortheilhaftesten ist es, die-
sen Spiegel so anzubringen, dass seine Ebene nahe senkrecht zu der Schneide ist, 
obwohl man darin nicht zu ängstlich zu sein braucht. In dem ersten der beiden 
Fälle bleibt der Spiegel, während des Spiels des Waagebalkens, sich selbst pa-
llel, im zweiten nicht; im ersten Falle wird also das
reflektirte Bild einer vor dem Spiegel aufgestellten Scale im Fernrohre unbeweglich 
sein, im zweiten wird es eine, wenn auch noch so kleine, Verrückung erleiden. 
Die Prüfung wird unendlich empfindlicher, wenn man das Spiegelbild nicht blos 
mit dem Auge, sondern mittelst eines Fernrohrs betrachtet, wobei zugleich jede 
äusserst kleine Verrückung des Bildes sicher und scharf mit einem Fernrohre er-
kennen.

Der Hofr.\ Gauss gebrauchte als Gegenstand eine etwa 5 Meter vor dem 
Spiegel vertical aufgerichtete, in Millimeter eingetheilte Scale; das 35 mal ver-
grössernde Fernrohr stand in nahe eben so grosser Entfernung. Es erschien so 
das Bild eines Millimeters etwa 20 Secunden gross, wovon man noch Zehntel 
schätzen kann. So lange die Schneide noch nicht vollkommen berichtigt war, 
ging das Bild der Scale an dem Fadenkreuze des Fernrohrs auf das regelmässigste 
auf und ab, wie der Waagebalken seine Schwingungen machte.

Für mathematisch gebildete Leser bedarf es blos der Andeutung, dass auf 
diese Weise nicht blos erkannt werden kann, nach welcher Seite eine Divergenz 
statt findet, sondern auch, hinreichend genau, wie gross dieselbe ist, wodurch, 
verbunden mit der Kenntniss der Weite der Gewinde der Correctionsschrauben, 
das Correctionsgeschäft in einen sichern Gang gebracht wird.

Der Vollständigkeit wegen mögen noch ein Paar andere Umstände hier er-
wähnt werden.

Wenn man einen etwas grossen Spiegel anwendet (der vom Hofr.\ Gauss 
gebrauchte, auf das Tragestück vermittelst einer eigenen Vorrichtung befestigte, 
hat 75 Millimeter Höhe), so ist es nothwendig, die Schalen mit hinlänglich schwe-
ren Gewichten zu belasten, weil sonst das Tragestück seitwärts umschlagen würde.

Es ist oben vorausgesetzt, dass die zu prüfende äussere Schneide mit der 
mittleren in Einer Ebene liege, also, wenn man die mittlere genau horizontal ge-
stellt hat, bei horizontalem Stande des Waagebalkens gleichfalls horizontal sei, 
und nur etwa seitwärts divergire. Gewöhnlich wird aber diese Voraussetzung auch 
nicht in äusserster Schärfe statt finden, sondern, die äussere Schneide bei jener 
Stellung etwas geneigt, oder das eine Ende etwas höher sein können als das an-
dere. Man erkennt dies, bei der beschriebenen Prüfungsmethode, daran, wenn 
beim Sinken des einen Hebelarmes das Bild in dem Fernrohre aufwärts, beim 
Steigen desselben aber abwärts geht, oder umgekehrt. In jenem Falle muss das 
nähere Ende der Schneide gehoben, in diesem gesenkt werden, wobei die schon 
erwähnten Correctionsschrauben gleichfalls das gehörige Mittel bieten. Die hier 
notwendige Correction ist übrigens von selbst so ziemlich erreicht, wenn man die 
frühere so weit getrieben hat, dass das Bild bei dem Spiele der Waage in der 
Höhe unverrückt bleibt.

\newpage
%==============================================================================
% SECTION: Physikalische Beobachtungen
% Pages 527-528
%==============================================================================
\section*{Physikalische Beobachtungen}
\addcontentsline{toc}{section}{Physikalische Beobachtungen}

\begin{center}
\fbox{\parbox{0.8\textwidth}{\centering\small
(Figürliche Darstellungen und Diagramme zu den physikalischen Beobachtungen,\\
wie im Original auf den Seiten 527--528.)
}}
\end{center}

\newpage
%==============================================================================
% SECTION: Magnetische Beobachtungen (Nordlicht)
% Pages 529-531
%==============================================================================
\section*{Magnetische Beobachtungen}
\addcontentsline{toc}{section}{Magnetische Beobachtungen}

\textit{Göttingische gelehrte Anzeigen. Stück 33. Seite 321 bis 324. 1831 Februar 28.}

Herr Prof.\ \textsc{Gerling} in Marburg hat der Königl.\ Societät eine Notiz über 
seine Wahrnehmung
\begin{center}
\textbf{des am 7. Januar d.\ J. gesehenen Nordlichts}
\end{center}
vorgelegt, welche zwar im Allgemeinen mit dem, was von andern Orten her bereits 
bekannt geworden ist, übereinstimmt, aber daneben noch einen, besonderer Auf-
merksamkeit werthen, und wie es scheint bisher noch nicht hinlänglich gewür-
digten Umstand berührt, daher wir hier einen Auszug aus derselben mittheilen.

Das Phänomen war in Marburg schon von 6 Uhr an gesehen. Herr \textsc{Ger}-
ling erhielt aber erst um 8 Uhr eine Benachrichtigung davon, und damals war 
am ganzen nördlichen Himmel, so tief herab wie die Aussicht aus den Fenstern 
seiner Wohnung reichte, gar nichts Ungewöhnliches zu erkennen. Allein gegen 
9 Uhr zeigten sich wieder auffallende rothe Streifen am nördlichen Himmel, und 
Herr \textsc{Gerling} begab sich sogleich auf den eine freie Aussicht beherrschenden 
Schlossberg, um noch so viel thunlich von der Erscheinung wahrzunehmen.

Zuerst wurden in einer Ausdehnung von etwa 50--60 Grad zwischen N.O. 
und N.W. blos rothe Streifen und Flecken am Himmel bemerklich, welche sich 
ohne vollständige Continuität in dem angegebenen Bogen im Azimuth und im 
Mittel etwa bis zu 45 Grad Höhe erstreckten. In der Mitte jenes Azimuthalbo-
gens um den Meridian herum und nach einer Schätzung etwa in 30--40 Grad 
Azimuthalausdehnung zeigten sich schwarze Flecke am sonst heitern Himmel, 
dem Ansehn nach mit nichts anderm als schwarzen Wölkchen zu vergleichen. 
Diese Flecke vermehrten sich allmählig, und bildeten endlich zusammenlaufend 
das dunkle Segment, welches nach allen Beschreibungen bei dem Nordlicht cha-
racteristisch zu sein scheint, indem zu gleicher Zeit die ersterwähnten rothen 
Flecken an Intensität zunahmen, und sich strahlenförmig gegen das schwarze 
Segment grupirten, von welchem aus zwischen den rothen Strahlen dann auch 
weisse und gelbliche erschienen, die ohne auffallend plötzliches Fortschiessen sich 
auf etwa 50 Grad in die Höhe erstrecken mochten.

So weit, fährt Herr \textsc{Gerling} fort, scheint diese Beobachtung mit dem, was 
andere Beobachter zu gleicher Zeit und bei früheren Nordlichtern gesehen haben, 
ganz übereinzustimmen, und würde also kaum eine Erwähnung verdienen, wenn 
nicht ein Umstand dabei mir aufgefallen wäre, welcher meines Wissens weder 
bei Gelegenheit dieses jetzigen Nordlichts, noch, so viel ich habe auffinden kön-
nen, sonst zur Sprache gekommen ist. Nemlich, nicht blos die Sterne des 
Schwans, über welchen die weissen und rothen Strahlen mit ihrer grossen In-
tensität hinweggingen, sondern auch der Stern $\alpha$ in der Leyer, welcher tief im 
schwarzen Segment stand, verloren an Sichtbarkeit und scheinbarer Helligkeit au-
genfällig gar nichts. Diese Thatsache scheint über die räthselhafte Frage, welche 
Bewandniss es mit dem dunkeln Segment eigentlich habe, wenigstens das negative 
Resultat zu geben, dass es keine gewöhnliche Wolke ist, weil solche für das Stern-
licht nicht permeabel sein könnte.

Schon bei dem Nordlicht vom 22.\ October 1804 bemerkte \textsc{Wrape}, allein 
ohne diesen Grund beizufügen, dass man das dunkle Segment unrichtig eine 
Wolke nenne, während \textsc{Gilbert} den Ausdruck in Schutz nimmt, und hinzusetzt, 
er habe im dunklen Segment nichts bemerkt, was ihn hätte auf den Gedanken 
bringen können, dass er dort etwas anderes als eine dunkele Wolke sähe. Auch 
die Meinung \textsc{Mayer}'s im Handbuch der physischen Astronomie, dass die dichtere 
mit Dünsten erfüllte Luft des Horizonts hinlänglich sei, das dunkle Segment zu 
erklären, scheint sich mit der von Herrn \textsc{Gerling} bemerkten Thatsache nicht ver-
einigen zu lassen.

Herr \textsc{Gerling} fügt noch bei, dass in den frühern Stunden, wo das in sei-
er Ausdehnung veränderliche Segment sich sehr hoch erstreckte, ein glaubwür-
diger Zeuge den Stern $\alpha$ Leyer in dem Segmente so hell wie zu irgend einer an-
dern Zeit glänzen gesehen, und ein anderer, zu einer Zeit, wo das dunkle Seg-
ment sich noch nicht bis zu jenem Sterne erstreckte, andere Sterne in dem Seg-
ment erblickt habe.

Herr \textsc{Gerling} hat noch einen Auszug aus seinem meteorologischen Journal 
vom 5.--9.\ Januar beigefügt, welcher jedoch ausser einem dreiviertel Zoll betra-
genden Steigen des Barometers vom 6.\ Januar Nachmittags bis 7.\ Januar Abends 
nichts auffallendes darbietet. Der Wind ging am 7.\ Januar aus Norden.

Die hier in Göttingen von Herrn Prof.\ \textsc{Hardung} an diesem Nordlichte ge-
machten Wahrnehmungen stimmen im Wesentlichen mit den von andern Orten 
bekannt gewordenen überein, doch verdient der Umstand erwähnt zu werden, 
dass während der Dauer des Phänomens die Magnetnadel um etwa dreiviertel 
Grad von ihrer gewöhnlichen Stellung nach Norden ging, und am andern Mor-
gen wieder auf dieselbe zurückgekommen war.

\textit{Göttingische gelehrte Anzeigen. Stück 128. Seite 1265 bis 1274. 1834 August 9.}

\begin{center}
\textbf{Magnetisches Observatorium}
\end{center}

Wir verdanken der Huld unserer Regierung ein neues, einem wichtigen 
Theile der Naturwissenschaften gewidmetes Institut, ein eignes 
\begin{center}
\textbf{für die magnetischen Beobachtungen und Messungen}\\
\textbf{errichtetes Observatorium.}
\end{center}

Obgleich der Bau desselben bereits im vorigen Herbst, und die innere Einrich-
tung seit Anfang dieses Jahres so weit vollendet ist, dass seit den ersten Monaten 
tägliche Beobachtungen angestellt werden konnten, so haben wir doch bisher An-
stand genommen, in diesen Blättern einen Bericht davon zu geben, weil wir erst 
einige Resultate der Beobachtungen damit verbinden zu können gewünscht haben. 
Die nach neuen Principien construirten magnetischen Apparate, welche im Jahre 
1832 in der hiesigen Sternwarte aufgestellt sind, haben wir bereits früher in die-
sen Blättern [Anzeige d.\ Intensitas v.\ m.] ausführlich beschrieben, und die damit 
erreichbare Schärfe ist aus dem dort Angeführten hinreichend ersichtlich: allein 
um diese Schärfe ganz zu erreichen, war eine Ausführung in grösserm Maassstabe, 
und um den Resultaten eine vollkommene Reinheit von fremden Einflüssen zu 
verschaffen, war ein besonderes eisenfreies Gebäude unumgänglich nöthig.

Das magnetische Observatorium, auf einem freien Platze, etwa hundert 
Schritt westlich von der Sternwarte errichtet, ist ein genau orientirtes längliches 
Viereck von 32 Par.\ Fuss Länge und 15 Fuss Breite, mit zwei Vorsprüngen an 
den längeren Seiten; der westliche Vorsprung bildet den Eingang, und dient zu-
gleich bei gewissen Beobachtungen als Erweiterung des Hauptsaals; der östliche 
Vorsprung, vom Hauptsaal ganz geschieden, dient zum Aufenthalt des Nacht-
wächters der Sternwarte. Im ganzen Gebäude ist ohne Ausnahme alles, wozu 
sonst Eisen verwandt wird, Schlösser, Thürangeln, Fensterbeschläge, Nägel u.\ s.\ w. 
von Kupfer. Für Abhaltung alles Luftzuges ist nach Möglichkeit gesorgt. Die 
Höhe des Saals ist etwas über 10 Fuss.

Der magnetische Apparat stimmt im Wesentlichen mit den oben erwähn-
ten überein, daher wir uns darauf einschränken, nur die Verschiedenheiten an-
zugeben. Der Magnetstab ist aus Uslarschem Gussstahl, welcher sich zu magne-
tischen Versuchen vortrefflich qualificirt; es wird von Zeit zu Zeit mit verschie-
denen Stäben gewechselt, die alle nahe gleiche Grösse haben, nemlich eine Länge 
von 610, Breite von 37, Dicke von 10 Millimetern; das Gewicht gegen vier Pfund. 
Der Spiegel ist 75 Millimeter breit und 50 hoch. Aufgehängt ist der Stab von 
der Mitte der Decke des Saals an einem 200fachen 7 Fuss langen ungedrehten 
Seidenfaden; der Torsionskreis ist aber nicht wie früher am obern Ende des Fa-
dens, sondern am untern, und mit dem Schiffchen, welches den Stab trägt, dreh-
bar verbunden. Seidene Aufhängungsfäden haben vor metallenen, wie bereits 
in der Abhandlung des Hofr.\ Gauss (\textit{Intensitas vis magneticae terrestris} Art.\ 9.) 
bemerkt ist, den grossen Vorzug, dass ihre Torsionskraft sehr klein ist; bei dem 
gegenwärtigen Tragfaden ist diese nur der Neunhundertste Theil der horizontalen 
Directionskraft des Magnetstabes, während die Torsionskraft eines Metallfadens 
von gleichem Tragvermögen etwa zehnmal stärker sein würde. Dagegen haben 
Seidenfäden, besonders wenn ihr Tragvermögen das an ihnen hangende Gewicht 
nicht weit übersteigt, die Inconvenienz, sich in den ersten Wochen, oder bei 
bedeutend verstärkter Belastung, beträchtlich zu verlängern; inzwischen wird 
dieser Inconvenienz hier durch den sinnreichen von Herrn Prof.\ \textsc{Weber} angege-
benen an der Decke befindlichen Aufhängungsapparat abgeholfen, womit der Fa-
den leicht, so viel nöthig, wieder aufgewunden werden kann, ohne seinen Platz 
zu verändern; zugleich aber kann dieser Apparat eben so leicht an der Decke 
verschoben werden, wenn im Lauf der Zeit die Veränderung der magnetischen 
Declination dies nöthig machen wird. Der Theodolith steht bisher auf einem sehr 
solide gearbeiteten hölzernen Stativ über einem besondern steinernen Fundament, 
und von dem Platze desselben ist durch das nördliche Fenster einer der Stadt-
thürme sichtbar, dessen Azimuth auf das genaueste bestimmt ist. Als Berichti-
gungsmarke für die unverrückte Stellung des Theodolithen dient blos ein zarter 
verticaler Strich an der gegenüberstehenden nördlichen Wand. Zum gewöhnli-
chen Gebrauch dient eine in Millimeter getheilte Scale von 4 Fuss Länge; für 
einige Beobachtungen wird dieselbe mit einer zwei Meter langen vertauscht. Der 
Werth eines Scalentheils ist $21{,}''3$. Für nächtliche Beobachtungen wurde bisher 
die Scale mit starker Beleuchtung versehen.

\textit{Göttingische gelehrte Anzeigen. Stück 128. Seite 1265 bis 1274. 1834 August 9.\ (Fortsetzung)}

Eine der Hauptanwendungen des Apparats besteht nun in der scharfen Be-
stimmung der magnetischen Declination und ihrer Veränderung in verschiedenen 
Tagesstunden, Monaten und Jahren. Alle Tage wird die Aufzeichnung zweimal 
zu bestimmten Stunden gemacht: man hat dazu die Vormittagsstunde 8 Uhr, und 
die Nachmittagsstunde 1 Uhr gewählt, mit welchen Zeiten bei regelmässigem 
Verlauf der täglichen Variationen die kleinste und die grösste Declination, we-
nigstens in den ersten Monaten des Jahrs, ungefähr zusammenfallen: dieser Auf-
zeichnung allemal genau bei derselben Uhrzeit hat man, aus wichtigen hier nicht 
weiter auszuführenden Gründen, vor dem jedesmaligen Abwarten des Minimum 
und Maximum unbedingt den Vorzug geben müssen. Diese Aufzeichnungen ha-
ben zwar schon seit dem 1.\ Januar den Anfang genommen: allein da zuerst 
noch mit den kleinern Apparaten gearbeitet wurde, die erst allmählig durch den 
gegenwärtigen ersetzt worden sind, so mag von den Resultaten erst künftig die 
Rede sein. Die Aufzeichnungen im März, April, Mai und Junius waren:

\begin{center}
\begin{tabular}{l@{\quad}l@{\quad}l}
\toprule
& 8 Uhr Vorm. & 1 Uhr Nachm.\\
\midrule
März, zweite Hälfte & $18^\circ\;38'\;16{,}''0$ & $18^\circ\;46'\;40{,}''4$ \\
April & \quad 36\quad 6{,}9 & \quad 47\quad 3{,}8 \\
Mai & \quad 36\quad 28{,}2 & \quad 47\quad 38{,}4 \\
Junius & \quad 37\quad 40{,}7 & \quad 47\quad 59{,}5 \\
Julius & \quad 38\quad 2 & \quad 48\quad 19{,}0 \\
\bottomrule
\end{tabular}
\end{center}

Ferner werden an gewissen bestimmten Tagen im Jahre 44 Stunden hin-
durch ununterbrochen in kurzen Zeitfristen die Veränderungen der Declination 
beobachtet. Man hat dazu dieselben bereits vor mehreren Jahren durch Herrn 
von \textsc{Humboldt} festgesetzten Tage gewählt, an welchen nach Verabredung schon 
an vielen zum Theil sehr entlegenen Plätzen ähnliche Aufzeichnungen mit \textsc{Gan}-
\textsc{sey}'schen Apparaten gemacht werden. Bis jetzt sind hier diese Beobachtungen 
dreimal angestellt, nemlich den 20.\ 21.\ März; 4.\ 5.\ Mai; 21.\ 22.\ Junius, und es 
haben daran Theil genommen ausser dem Hofr.\ Gauss die Herren Prof.\ \textsc{Weber}, 
Prof.\ \textsc{Ulrich}, Dr.\ \textsc{Weber}, Dr.\ \textsc{Goldschmidt}, Dr.\ \textsc{Listing}, \textsc{Sartorius}, \textsc{Denzler} und 
\textsc{Wilh.\ Gauss}. Der Zweck dieser Beobachtungen ist, theils den regelmässigen 
Verlauf nach und nach immer vollständiger kennen zu lernen, theils die Be-
wandtniss, welche es mit den so häufig dazwischen kommenden, zuweilen, beson-
ders bei Nordlichtern, ungemein beträchtlichen ausserordentlichen Anomalien 
hat, durch Vergleichung der gleichzeitigen Beobachtungen an verschiedenen Or-
ten zu erforschen. Die Aufzeichnungen geschahen hier, im März von 20 zu 
20 Minuten, und zum Theil in halb so grossen Zwischenzeiten; im Mai von 10 
zu 10 Minuten und zum Theil in doppelt engen Grenzen; im Junius durchge-
hends von 5 zu 5 Minuten. Anomalien wurden hier bemerkt, ein Paar auffal-
lend grosse in der Nacht vom 20.\ zum 21.\ März; sehr bedeutende und zahlreiche 
in den Nächten vom 4.\ und 5.\ Mai; und einige zwar nicht grosse aber doch be-
stimmt hervortretende am 21.\ Junius, während den ganzen 22.\ Junius der Ver-
lauf überaus regelmässig war. Von denjenigen correspondirenden Beobachtun-
gen, welche, wie schon erwähnt, Herrn von \textsc{Humboldt} ihre Veranlassung ver-
danken, sind uns bisher keine bekannt geworden, als die vom 4.\ und 5.\ Mai 
an zwei entlegenen Orten, in Christiania und in Dorpat. Dagegen sind mit 
der hiesige construirten Apparate die correspondirenden Beobachtungen vom 4. 
und 5.\ Mai auf einem Gute in Baiern, einige Meilen südlich von Meiningen, sehr 
vollständig angestellt, woraus eine wahrhaft bewundernswürdige Übereinstim-
mung mit den hier beobachteten grossen Anomalien nach Zeit, Grösse und Wech-
sel derselben hervorgeht, so dass man in den graphischen Darstellungen die eine 
beinahe als eine Copie der andern mit allen barocken durch jene Anomalien her-
vortretenden Figuren ansehen möchte. Ein eben so schöner Erfolg hat sich am 
21.\ und 22.\ Junius gezeigt, wo correspondirende Beobachtungen in Berlin zum 
ersten Male mit einem dem hiesigen ähnlichen obwohl kleinern Apparate von 
Herrn Prof.\ \textsc{Encke} unter Beistand von Herrn \textsc{Poggendorff}, \textsc{Mädler} und \textsc{Wolfers} 
angestellt wurden. Auch dort waren keine andere Anomalien, als die hier 
beobachteten, aber diese fast treu copirt, und eben dasselbe zeigten die von Herrn 
\textsc{Sartorius} damals in Frankfurt am Main gemachten Beobachtungen. Diese Re-
sultate können bereits als eine schöne Frucht der verabredeten Beobachtungen 
angesehen werden, da daraus auf das klarste hervorgeht, dass kleinere und 
grössere Anomalien der Magnetnadel, die zuweilen in ziemlich kurzen Fristen 
wechseln, nicht locale, sondern kräftige, weithin wirkende Ursachen haben müs-
sen, was man in Beziehung auf sehr grosse mit Nordlichtern in Verbindung ste-
hende Unregelmässigkeiten auch schon früher bemerkt hatte. So wie in Zu-
kunft die Theilnahme an diesen verabredeten Beobachtungen mit den eben so 
scharfen als bequemen Apparaten sich immer weiter ausbreiten wird, wozu schon 
die schönsten Aussichten vorhanden sind, wird es nicht fehlen, dass wir über 
diese höchst merkwürdigen und räthselhaften Erscheinungen umfassende Auf-
klärung erhalten.

\medskip

Drei Bestimmungen mit verschiedenen Stäben gaben
\[
= 1{,}792;\quad = 1{,}783;\quad = 1{,}786
\]
als Werth der horizontalen Kraft, wobei, wie bei den frühern Bestimmungen 
mit kleinern Stäben, deren geringe Verschiedenheit von den gegenwärtigen man 
mit Vergnügen bemerken wird, die Zeitsecunde, das Millimeter und das Milli-
gramm als Einheiten zum Grunde liegen.

Eben so, wie mit dem frühern in der Sternwarte aufgestellten Apparate, 
hat man nun auch mit dem gegenwärgen im magnetischen Observatorium Vor-
richtungen zu electromagnetischen Versuchen und Messungen verbunden. Der 
aufgehängte Magnetstab ist von einem aus 200 Umwindungen bestehenden Mul-
tiplicator umgeben, dessen Construction die Anwendung von nichtbesponnenem 
Draht erlaubte: die Drahtlänge beträgt 1100 Fuss. Mit Hülfe eines sehr ein-
fach construirten Commutators kann der Beobachter, ohne sein Auge vom Fern-
rohr zu entfernen, jeden Augenblick die Richtung des galvanischen Stroms um-
kehren, oder den Strom ganz unterbrechen.

Wir können hiebei eine mit den beschriebenen Einrichtungen in genauer 
Verbindung stehende grossartige und bisher in ihrer Art einzige Anlage nicht 
unerwähnt lassen, die wir unserm Herrn Prof.\ \textsc{Weber} verdanken. Dieser hatte 
bereits im vorigen Jahre von dem physikalischen Cabinet aus über die Häuser 
der Stadt hin bis zur Sternwarte eine doppelte Drahtleitung hergestellt, deren 
fortgesetzter Gebrauch bei galvanischen Versuchen die vollkommenste Beständig-
keit gezeigt hat; inzwischen hat er sie in diesem Jahre nicht nur durch eine 
zweite, sondern auch durch eine öffentlich frei in der Luft über die Dächer hin-
geführte von derselben Länge vermehrt. Jede dieser vier Leitungen hat eine 
Länge von mehr als 1500 Fuss, so dass mit ihnen jedes beliebige der drei In-
stitute, die sie verbinden, entweder für sich oder mit einem oder beiden andern 
in galvanische Verbindung gebracht werden kann.

\medskip

Eben so interessant, wie die rein magnetischen Beobachtungen sind die mit 
diesem Apparat anzustellenden electrodynamischen Versuche. Zu diesem Zweck 
ist der Stab von einem ähnlichen Multiplicator umgeben, wie der Stab des mag-
netischen Observatoriums, nur dass jener grössere Dimensionen, und eine Draht-
länge von 2700 Fuss in 270 Umwindungen hat. Dieser Multiplicator ist in die 
grosse schon in dem frühern Bericht erwähnte Drahtkette gebracht, welche die 
Sternwarte, das magnetische Observatorium und das physicalische Cabinet ver-
bindet, und in welcher der galvanische Strom zusammen eine Drahtlänge von 
11000 Fuss, also fast einer halben geographischen Meile zu durchlaufen hat, 
und dann drei magnetische Apparate zugleich afficirt, nemlich
\begin{enumerate}
\item[I.] den 25pfündigen Stab in der Sternwarte,
\item[II.] den 4pfündigen Stab im magnetischen Observatorium\\(Multiplicator von 200 Umwindungen)
\item[III.] den einpfündigen Stab im physikalischen Cabinet\\(Multiplicator von 160 Umwindungen).
\end{enumerate}
Einzelne Theile der Kette können in vielfachen Combinationen nach Gefallen 
mit Leichtigkeit abgesperrt werden.

Von den zahlreichen Versuchen, welche schon jetzt mit diesen Apparaten 
gemacht sind, führen wir hier nur einige an.

Wenn ein galvanischer Strom mit der Kette in Verbindung gesetzt wird, 
so erscheinen die Bewegungen der Magnetstäbe in den drei Apparaten so augen-
blicklich, dass ihr Anfang sich auf einen kleinen Bruch einer Zeitsecunde genau 
beobachten lässt. Die Vergleichung der Uhren bei den drei Apparaten liefert so 
vollkommen übereinstimmende Resultate, der Strom möge an dem einen Ende, 
oder an dem andern, oder in der Mitte erzeugt sein, dass daraus die Unmessbar-
keit der Zeit, in welcher der Strom eine halbe Meile durchläuft, vollkommen be-
stätigt wird. Nach den interessanten Versuchen von \textsc{Wheatstone}, welche neuer-
lich in den Philosophical Transactions für 1834 bekannt gemacht sind, und nach 
welchen der electrische Strom im Metall eine gewisse Zeit braucht, möchte es 
scheinen, dass bei einer so grossen Ausdehnung der Kette die Zeit, welche der 
Strom nöthig hat, um sich von dem einen Ende zum andern zu verbreiten, nicht 
mehr ganz unmerklich sein würde. Allein bei unsern Versuchen zeigt sich da-
von nicht die mindeste Spur, und da der von uns bemerkte Umstand die Mög-
lichkeit darbietet, diese Zeit, wenn sie gleich eine Million Mal kürzer wäre, als 
nach \textsc{Wheatstone}'s Versuchen, zu bemerken, so können wir nur schliessen, dass 
bei metallischen Leitungen von der Länge einer halben Meile die Fortpflanzung 
der Elektricität eine ganz unendlich kleine Zeit erfordert.

Die Ablenkungen, welche derselbe Strom an den drei Apparaten hervor-
bringt, sind freilich nicht unmittelbar mit einander vergleichbar, da sie sich 
offenbar in den drei Apparaten mit verschiedenen Einheiten, welche von den Di-
mensionen der Multiplicatoren und der Geltung der Scalentheile in Bogensecun-
den abhängen. Nun zeigen aber zahlreiche angestellte Versuche, dass zwischen 
den Ablenkungen an den drei Apparaten durch denselben Strom in einerlei Au-
genblick stets genau ein constantes Verhältniss Statt findet, der Strom möge an 
dem einen, oder an dem andern Ende, oder in der Mitte erzeugt sein. Es er-
gibt sich daraus das wichtige Resultat, dass der Strom in seiner ganzen Länge 
dieselbe Intensität hat, wenigstens nichts merkliches davon verliert. Man wird 
in Zukunft besonders aufmerksam darauf sein, ob dieses Resultat auch unter ei-
genthümlichen Umständen, namentlich während starken Regens, seine Gültig-
keit behält.

Bei allen drei Apparaten sind Commutatoren (Gyrotrope) mit der Kette 
verbunden, wodurch man die Richtung des Stroms mit Leichtigkeit umkehren 
kann. Dem Commutator in der Sternwarte hat der Hofr.\ Gauss eine eigenthüm-
liche Einrichtung gegeben, wonach diese Umkehrung durch einen einzigen Druck 
mit dem Finger, also ganz augenblicklich, bewirkt wird. Wenn man diese Um-
kehrung, immer in so grossen Zeitfristen wie die Schwingungsdauer des Einen 
Stabes, wiederholt ausführt, so werden die Schwingungen dieses Stabes immer 
grösser. Man hat dies zu einem Experiment benutzt, wobei eine auffallende 
mechanische Wirkung hervorgebracht wird. Herr Prof.\ \textsc{Weber} liess zur Seite 
des Magnetstabes im physikalischen Cabinet eine leichte Auslösung für einen 
Wecker oder eine Pendeluhr anbringen. Dieses Auslösen gelingt jedesmal durch 
den von der Sternwarte aus geleiteten Strom nach ein Paar Schwingungen auf das 
vollkommenste. Dass man mit dem 25pfändigen Stabe eine noch viel stärkere 
mechanische Wirkung würde hervorbringen können, leuchtet von selbst ein.

Besonders wichtige Dienste leisten diese Apparate bei der Erforschung der 
mathematischen Gesetze, nach welchen die Induction eines magnetischen Ele-
ments auf die galvanische Kette Statt findet, wofern das Element sich der Kette 
nähert oder von ihr entfernt. Man bedient sich hiezu einer Rolle von blankem 
Kupferdraht, die auf einem Mittelstück aus Messing aufgewunden ist, und die 
man auf den Magnetstab oder über denselben hin aufschiebt, oder davon abzieht. 
Diese Rolle kann über die freistehende Hälfte eines starken Magnetstabes ge-
führt werden, und während dieser Operation geht allemal durch die Kette ein 
galvanischer Strom, ein starker, aber von kurzer Dauer, oder ein schwächerer 
von längerer Dauer, je nachdem die Manipulation schneller oder langsamer ge-
schieht, so dass die Gesammtwirkung Eines Aufschiebens von der Schnelligkeit 
der Operation unabhängig ist. Der Strom an sich dauert immer nur so lange, 
wie die Bewegung der Rolle. Das Abziehen der Rolle bringt einen entgegenge-
setzten Strom hervor, eben so das Aufschieben mit dem entgegengesetzten Ende. 
Geschieht die Bewegung sehr schnell, so ist die Wirkung des Stroms auf die Mag-
netnadel in einem der mit der Kette verbundenen Multiplicatoren einem augen-
blicklichen Stosse von bestimmter Stärke gleich zu setzen. Abziehen und ver-
kehrt wieder Aufstecken bewirkt also zwei gleichnamige Impulse der Magnetna-
del, und ein neues Abziehen und wieder umgekehrt Aufschieben würde daher 
zwei unter sich gleiche, aber den vorigen entgegengesetzte Impulse hervorbrin-
gen; allein wenn dazwischen der Commutator gewechselt ist, so geschehen auch 
die letzten beiden Wirkungen in demselben Sinn, wie die beiden ersten. Ein 
solcher vollständiger Wechsel (Abziehen, Verkehrtaufstecken und Commutator-
umstellung) geschieht ganz bequem in zwei Secunden, und man kann daher, wenn 
man will, während einer Schwingungsdauer des grossen Magnetstabes bequem 
und tactmässig 21 Wechsel vollenden, und dadurch letztern in so starke Bewe-
gung bringen, dass die ganze Scale aus dem Gesichtsfelde des Fernrohrs geht. 
Diese Andeutung wird hinreichen zu übersehen, wie die Stärke des durch diese 
Inductionsart entstehenden galvanischen Stroms mit Schärfe gemessen werden 
kann. Diese Stärke hängt aber zugleich von dem Widerstande ab, welchen die 
Kette selbst darbietet, und nimmt mehr oder weniger zu, je nachdem mehr oder 
weniger Stücke der Kette abgesperrt sind.

Wir haben oben erwähnt, dass die Abnahme des Schwingungsbogens bei 
der grossen Nadel in verschiedenen Zeiten sehr ungleich gewesen ist. Ähnliche 
Verschiedenheiten hatten sich schon im Jahr 1832 bei den kleinen Apparaten ge-
zeigt, auch später bei der Nadel im magnetischen Observatorium: allein diese 
Verschiedenheiten bleiben immer innerhalb viel engerer Grenzen, als bei dem 
Stabe der Sternwarte, wo die Abnahme des Schwingungsbogens von einer Schwin-
gung zur folgenden in verschiedenen Versuchsreihen zwischen $\tfrac{1}{650}$ und $\tfrac{1}{5}$ 
schwankte. Diese merkwürdige Erscheinung hat die Aufmerksamkeit des Hofr.\ 
Gauss besonders auf sich gezogen, und es scheint dabei ein Zusammentreffen meh-
rerer Ursachen Statt zu finden, die zum Theil noch jetzt räthselhaft bleiben: in-
zwischen ist es dem Hofr.\ Gauss gelungen, diejenige Ursache, welche bei weitem 
den stärksten Einfluss hat, auszumitteln. Er bemerkte nemlich, dass allemal 
der Schwingungsbogen viel schneller abnahm, wenn die Kette geschlossen, als 
wenn sie offen war, und so war es leicht, als Ursache jener schnellen Abnahme, 
die Reaction eines in der Kette durch die Schwingung der Nadel selbst, vermöge 
der Induction, erzeugten galvanischen Stroms zu erkennen, welcher bei der fol-
genden Rückschwingung die entgegengesetzte Richtung hat, und stets auf Ver-
minderung des Schwingungsbogens wirkt.

\newpage
%==============================================================================
% SECTION: Beobachtungen der magnetischen Variation
% Pages 537-540
%==============================================================================
\section*{Beobachtungen der magnetischen Variation in Göttingen und Leipzig}
\addcontentsline{toc}{section}{Beobachtungen der magnetischen Variation}

\textit{Göttingische gelehrte Anzeigen. Stück 36. Seite 345 bis 357. 1835 März 7.}

In der Sitzung der Königl.\ Societät am 14.\ Februar stattete der Hofr.\ Gauss 
einen Bericht über die in dem magnetischen Observatorium, und in Verbindung 
damit anderwärts gemachten Beobachtungen ab, woraus wir hier einen Auszug 
mittheilen, der als
\begin{center}
\textbf{eine Fortsetzung der am 9.\ August 1834 gegebenen Nachricht}
\end{center}
betrachtet werden kann.

Die täglich zweimaligen Aufzeichnungen des Standes der Nadel sind un-
unterbrochen fortgesetzt, und umfassen nun bereits beinahe ein volles Jahr. Die 
monatlichen Mittel, seit Julius v.\ J., waren:

\begin{center}
\begin{tabular}{l@{\quad}l@{\quad}l}
\toprule
& 8 Uhr Vorm. & 1 Uhr Nachm.\\
\midrule
1834 August & $18^\circ\;38'\;48{,}''1$ & $18^\circ\;49'\;11{,}''0$ \\
September & \quad 36\quad 58{,}4 & \quad 46\quad 32{,}3 \\
October & \quad 37\quad 18{,}4 & \quad 44\quad 47{,}2 \\
November & \quad 37\quad 27{,}4 & \quad 43\quad 4{,}3 \\
December & \quad 37\quad 54{,}8 & \quad 41\quad 32{,}7 \\
1835 Januar & \quad 37\quad 51{,}5 & \quad 42\quad 14{,}4 \\
\bottomrule
\end{tabular}
\end{center}

Die verabredeten Beobachtungen an bestimmten Tagen in kurzen ununter-
brochenen Zeitfristen, mit deren Einrichtung in den letzten Monaten einige an 
einem andern Orte bekannt gemachte Abänderungen getroffen sind, haben seit 
der letzten Nachricht an vier Hauptterminen Statt gefunden, einige ausserordent-
liche Nebentermine ungerechnet. Die Theilnahme an denselben hat sich bereits 
weiter ausgebreitet, und wird bald noch weiter verbreitet werden, auch sind dar-
aus schon sehr merkwürdige Resultate hervorgegangen, denen ähnlich, welche 
in dem frühern Bericht erwähnt wurden. Eine graphische Darstellung der Har-
monie unter den Beobachtungen vom 1.\ und 2.\ October, und vom 29.\ und 30.\ 
November in Göttingen, Leipzig und Berlin, wird nächstens in \textsc{Poggendorff}'s 
Annalen der Physik erscheinen: noch merkwürdiger aber ist die Übereinstim-
mung der Beobachtungen vom 5.\ und 6.\ November in Copenhagen und Mailand 
in allen zahlreichen und auffallend grossen Schwankungen, von welchen gleich-
falls eine Zeichnung an einem andern Orte gegeben werden wird. Wir treten 
hier in eine Welt von geheimnissvollen Naturkräften, deren wunderbar wechseln-
des Spiel sich über den halben Durchschnitt von Europa, in gleichem Augen-
blick, und bis in die kleinsten Nuancen auf gleiche Weise, offenbart, und deren 
Wirkungskreis zu ermessen diese Standlinie noch viel zu klein erscheint.

\newpage
%==============================================================================
% SECTION: Beobachtungen vom 1. und 2. October 1834
% Pages 537-540 (continued)
%==============================================================================
\section*{Beobachtungen der magnetischen Variation in Göttingen und Leipzig am 1.\ und 2.\ October 1834}
\addcontentsline{toc}{section}{Beobachtungen vom 1.\ und 2.\ October 1834}

\textit{\textsc{Poggendorff}. Annalen der Physik und Chemie. 1834. Bd.\ 33.}

Die in meinem Aufsatze über das hiesige magnetische Observatorium er-
wähnten Beobachtungen der magnetischen Variation an den verabredeten Tagen 
sind seitdem hier noch zwei Mal angestellt, am 6.\ und 7.\ August, und am 23.\ 
und 24.\ September. Im ersten Termin kamen recht starke und merkwürdige 
Anomalien vor, und es ist daher um so mehr zu bedauern, dass zufällige Ursa-
chen die Anstellung correspondirender Beobachtungen an andern Orten gehin-
dert haben. Die September-Beobachtungen sind hingegen ganz vollständig auch 
in Leipzig und Berlin und beinahe vollständig in Braunschweig angestellt; ausser-
dem auch zur Hälfte in Copenhagen, wo durch Versehen der 24.\ und 25.\ Sep-
tember anstatt des 23.\ und 24.\ genommen wurden. Die vollständige Bekannt-
machung dieser Beobachtungen würde jedoch geringeres Interesse haben, da der 
Verlauf an diesen beiden Tagen sehr regelmässig war, obgleich mehrere an sich 
sehr kleine Anomalien in den ersten 24 Stunden an allen vier Plätzen eine be-
wunderungswürdige Harmonie gezeigt haben. Merkwürdig bleibt indessen, dass, 
einer Zeitungsnachricht zufolge, am 23.\ September Abends in Glasgow ein sehr 
starkes Nordlicht gesehen worden ist, welches mithin ganz entschieden, wenig-
stens keinen sich bis Norddeutschland erstreckenden Einfluss auf die Magnetna-
del gehabt hat.

Die Anwesenheit des Herrn Prof.\ \textsc{Weber} in Leipzig veranlasste inzwischen, 
noch einige ausserordentliche Stunden zu gleichzeitigen Beobachtungen in Göt-
tingen und Leipzig festzusetzen, wozu die Tage 1.\ und 2.\ October Morgens $7\tfrac{1}{2}$ 
bis $8\tfrac{1}{2}$, Mittags $12\tfrac{1}{2}$ bis 14, und Abends 8 bis 10 Uhr gewählt wurden. Abge-
sehen von einigen kleinen Versäumnissen wurden diese Stunden an beiden Orten 
inne gehalten; im hiesigen magnetischen Observatorium beobachtete mein Sohn, 
\textsc{Wilhelm Gauss}, in Leipzig Herr Prof.\ \textsc{Weber}, Herr Prof.\ \textsc{Möbius} und Herr Dr.\ 
\textsc{Tellkamp}. --- Man wird nicht ohne Vergnügen die grosse Übereinstimmung 
nicht bloss in den grossen Bewegungen, welche am Abend des 1.\ October statt-
fanden, sondern fast in sämmtlichen kleinen bemerken, so dass deren Quellen 
sich als auf grosse Ferne hinwirkende, obwohl zur Zeit noch sehr räthselhafte 
Kräfte, auf das Unverkennbarste ausweisen. In Leipzig waren die Anomalien 
im Allgemeinen etwas kleiner als in Göttingen; letzterem Orte wird daher der 
Heerd der wirkenden Kräfte näher gewesen sein. Ich bemerke nur noch, dass 
während eines Theils jener Stunden ich selbst an einem zweiten in der hiesigen 
Sternwarte aufgestellten Apparat, wovon ich bald eine ausführlichere Nachricht 
zu geben gedenke, beobachtet habe, und dass diese Beobachtungen einen fast 
vollkommenen Parallelismus mit denen des hiesigen magnetischen Observato-
riums in den grösseren und kleineren Bewegungen ergeben haben; ein ähnlicher 
Erfolg hatte auch am 23.\ und 24.\ September, so wie bei vielen sonstigen Ver-
suchen, statt, in dem Maasse, dass schon öfters die Uhren an beiden Plätzen 
bloss mittelst der magnetischen Erscheinungen bis auf einen kleinen Bruchtheil 
einer Zeitminute genau verglichen werden konnten. Dasselbe gelingt mittelst 
der grösseren Bewegungen am 1.\ und 2.\ October zwischen Göttingen und Leip-
zig, wo an beiden Orten die Uhren nur wenige Secunden von der mittleren 
Ortszeit abwichen.

Durch diese Erfahrungen erhalten nun auch die kleinen, in sehr kurzen 
Zeitfristen wechselnden Schwankungen der Magnetnadel ein überaus grosses In-
teresse; man muss wünschen, dass auch diese durch die Beobachtungen an vie-
len von einander entfernten Plätzen sorgfältig verfolgt werden, und es wird da-
her unumgänglich nöthig alle Beobachtungen in recht kurzen Zeitintervallen zu 
machen. Bisher beobachteten wir von 5 zu 5 Minuten: aber auch dieses Inter-
vall ist noch fast zu lang, und wir denken künftig immer von 3 zu 3 Minuten 
den Stand der Magnetnadel an den verabredeten Tagen zu bestimmen. Ich darf 
dabei nicht unbemerkt lassen, dass das Verfahren, welches der Herr Herausge-
ber dieser Annalen (Bd.\ XXXI.\ S.\ 569 bis 572) erklärt hat, von uns nur an-
fänglich gebraucht, aber schon lange mit einem etwas abgeänderten vertauscht 
ist. Um den Stand der Magnetnadel für einen Augenblick zu erhalten, beob-
achten wir sie in sechs verschiedenen, immer um Eine Schwingungsdauer ge-
trennten Momenten, und so, dass der gewünschte Moment in die Mitte fällt. 
Anstatt der genauen Schwingungsdauer wird die nächste runde Zahl von Se-
cunden (oder vielmehr von Uhrschlägen) gewählt, z.\ B. im magnetischen Observato-
rium $20''$ anstatt $20{,}''$... Die Beobachtungen am 2.\ October für $8^h\,15'$ Abends 
standen daher so:

\begin{center}
\begin{tabular}{r@{\quad}l}
$8^h\;14'\;10''$ & 672{,}6 \\
\quad 30 & 672{,}3 \\
\quad 50 & 671{,}3 \\
$8^h\;15'\;10''$ & 671{,}8 \\
\quad 30 & 669{,}9 \\
\quad 50 & 670{,}8 \\
\end{tabular}
\end{center}

Hieraus ergeben sich fünf Mittel, die eigentlich den beigesetzten Zeiten cor-
respondiren:
\begin{center}
\begin{tabular}{r@{\quad}l}
$8^h\;14'\;20''$ & 672{,}45 \\
\quad 40 & 671{,}80 \\
$8^h\;15'\;0''$ & 671{,}55 \\
\quad 20 & 670{,}85 \\
\quad 40 & 670{,}35 \\
\end{tabular}
\end{center}

und daraus das Mittel für $8^h\;15'\;\;671{,}40$.

Ich habe absichtlich dieses Beispiel gewählt, wo die Nadel schnelle Ver-
änderungen zeigte, die selbst von 20 zu 20 Secunden sich so entschieden dar-
stellen. Wir haben Fälle genug, wo ein ähnlicher Erfolg selbst in halb so grossen 
Zeitintervallen eintritt. Gewisse Abänderungen in jener Beobachtungsart (die 
wir öfters anwenden) zu erklären, so wie die Rechtfertigung jener Art das Mit-
tel zu nehmen, die mit gutem Vorbedacht gewählt ist, muss ich mir für eine an-
dere Gelegenheit vorbehalten. Aber unerwähnt lassen darf ich nicht (da der 
Herr Herausgeber dieser Annalen a.a.O.\ es nicht ausdrücklich bemerkt hat), dass 
es eine wesentliche Bedingung für die Zulässigkeit aller dieser Beobachtungsar-
ten ist, die Nadel vorher so viel wie möglich beruhigt zu haben, so dass die 
Schwingungen nur eine geringe Anzahl von Scalentheilen betragen. Im hiesigen 
magnetischen Observatorium ist eine solche Beruhigung oder wenigstens eine 
Wiederholung derselben, im Laufe der Beobachtung selten nöthig. Wer aber 
in einem weniger günstigen Local beobachtet, darf durchaus nicht unterlassen, 
dies, so oft es nöthig wird, in der Zwischenzeit mit den bekannten Mitteln zu thun.

Da bei der gegenwärtig als nothwendig sich zeigenden Verengerung der 
Zwischenzeiten die Beobachtungen sehr viel mühsamer werden als früher, wo die 
Forderung sich auf die Aufzeichnung von Stunde zu Stunde beschränkte, so ist 
mehrseitig der Wunsch geäussert, künftig sowohl die Anzahl als die Dauer der 
Termine etwas zu verkürzen. ---

\begin{flushright}
Göttingen, den 5.\ November 1834.
\end{flushright}

\newpage
%==============================================================================
% SECTION: Weitere magnetische Beobachtungen
% Pages 541-548
%==============================================================================
\section*{Weitere magnetische Beobachtungen}
\addcontentsline{toc}{section}{Weitere magnetische Beobachtungen}

\textit{Göttingische gelehrte Anzeigen. Stück 128. Seite 1265 bis 1274. 1834 August 9.\ (Fortsetzung)}

Die hiesigen Einrichtungen für magnetische Beobachtungen haben inzwi-
schen mehrere wesentliche Erweiterungen erhalten. Für manche Beobachtungen 
ist, wenn grosse Schärfe verlangt wird, die Zuziehung eines zweiten Apparats, 
in einiger Entfernung vom Hauptapparate, unumgänglich nothwendig, um von 
den stündlichen Veränderungen der magnetischen Kraft Rechnung tragen zu kön-
nen. Zu diesem Zweck ist seit August v.\ J., nachdem die im Jahre 1832 ge-
brauchten Apparate an das physicalische Cabinet abgegeben sind, in der Stern-
warte ein grosser Magnetstab aufgehängt, mit übrigens ganz ähnlichem Zubehör, 
wie der Stab im magnetischen Observatorium. Der Magnetstab in der Stern-
warte, gleichfalls aus Uslarschem Gussstahl, ist 4 Fuss lang, fast drei Zoll breit 
und über einen halben Zoll dick, und wiegt 25 Pfund. Er hängt an einem 16 
Fuss langen tausendfachen Seidenfaden.

---

\textit{Göttingische gelehrte Anzeigen. Stück 36. Seite 345 bis 357. 1835 März 7.\ (Fortsetzung)}

Nachdem diese Erklärung gefunden war, war es leicht, den Erfolg eini-
ger Versuche vorauszusehen, welche wohl zu den auffallendsten im Gebiet des 
Electromagnetismus gerechnet werden dürfen, und selbst die quantitativen Ver-
hältnisse der Erscheinungen im Voraus zu berechnen, welche auch bei den wie-
derholt angestellten Versuchen stets auf das vollkommenste bestätigt sind. Es 
sind folgende.

Wenn der Magnetstab in der Sternwarte (I) in Schwingungen gesetzt wird, 
etwa so grosse wie der Kasten verstattet, so haben diese gar keinen Einfluss auf 
die Nadeln im magnetischen Observatorium (II) oder im physikalischen Cabinet 
(III), sondern diese bleiben in Ruhe, wenn sie vorher in Ruhe waren, vorausge-
setzt, dass die Kette offen, oder wenigstens die die letzten Nadeln einschliessen-
den Multiplicatoren davon abgesperrt sind. Allein in dem Augenblick, wo die 
Kette geschlossen oder z.\ B. der Multiplicator von II in die geschlossene Kette 
hineingebracht wird, fängt die Nadel II sogleich an mitzuschwingen. Ist die Na-
del II schon vorher in Schwingung gewesen, so erhalten die Schwingungen den 
eigenthümlichen Character gemischter Schwingungen, wovon die eine von dem 
Initialzustande abhängt, und dieselbe Periode hat, wie die Schwingungen dieser 
Nadel unter dem blossen Einfluss des Erdmagnetismus (20"), während die andere 
eine Periode von 42" befolgt (wie die grosse Nadel I), und ihre Grösse dem 
Schwingungsbogen von I proportional ist (etwa $\tfrac{1}{24}$, wenn die Kette hinter dem 
Multiplicator von II abgesperrt ist). Dies ist vollkommen mit den Resultaten 
der Theorie in Übereinstimmung, eben so wie der stets genau bestätigte Umstand, 
dass die Schwingungen von I und die inducirten Schwingungen von II, obwohl 
Perioden von gleicher Dauer, doch nicht gleichen Anfang haben, sondern stets 
eine halbe Schwingungszeit (21") in dieser Beziehung differiren, und zwar in dem 
Sinn, wie es nach den statt findenden Umständen die Theorie vorausbestimmt.

Was hier beispielsweise von der Nadel II gesagt ist, findet auf ganz ähnliche 
Weise bei der Nadel III Statt, deren natürliche Schwingungsdauer 14" beträgt, 
und die unter der Einwirkung der Induction zusammengesetzte Schwingungen 
von 14" und 42" Periode befolgt.

Ein ganz anderer Erfolg muss der Theorie zufolge in dem Fall Statt finden, 
wenn eine zweite Nadel, deren natürliche Schwingungsdauer genau eben so gross 
ist, wie die des grossen Magnetstabes, mit einem Multiplicator sich in der Kette 
befindet, in welcher der grosse Stab schwingt. Jene, so lange vollkommen ruhig, 
als die Kette offen ist, fängt gleichfalls in dem Augenblick an mitzuschwingen 
wo die Kette geschlossen wird, allein diese Schwingungen, von derselben Dauer, 
wie die natürlichen, nehmen an Grösse beständig zu, bis diese (erst nach sehr 
langer Zeit) zu einem Maximum kommt, wo der Widerstand der Luft der Ver-
grösserung durch die Inductionskraft das Gleichgewicht hält. Um diesen merk-
würdigen Versuch wirklich anstellen zu können, wurde (da die Aufhängung ei-
nes grossen Stabes wegen Mangel eines zweiten dafür passenden Multiplicators 
jetzt nicht thunlich war) der einpfündige Stab des physikalischen Cabinets durch 
Verbindung mit einem ähnlichen etwas schwächer magnetisirten auf bekannte 
Weise astatisch gemacht, oder vielmehr zu einer Doppelnadel, deren natürliche 
Schwingungsdauer genau auf 42{,}3 gebracht wurde. Der Versuch gelang damit 
auf das vollkommenste. Der in der Sternwarte schwingende Stab theilte dieser 
Doppelnadel im physicalischen Cabinet, in dem Augenblick wo die Kette ge-
schlossen wurde, wie durch eine wunderbare Sympathie seine Schwingungen mit, 
und zwar so, dass jede folgende etwa 50 Scalentheile oder einen halben Grad 
grösser wurde, als die vorhergehende. Bald ging das ganze Scalenbild aus dem 
Felde, allein fortwährend konnte man an der immer wachsenden Schnelligkeit, 
mit welcher das Scalenbild durch das Gesichtsfeld ging, die Zunahme des Schwin-
gungsbogens erkennen. Über eine Stunde wurde dies wunderbare sympatheti-
sche Spiel beobachtet.

Es braucht kaum bemerkt zu werden, dass auch der vierpfündige Stab im 
magnetischen Observatorium in die geschlossene Kette einen Strom inducirt, des-
sen Dasein an der schnellen Abnahme des Schwingungsbogens auf das bestimm-
teste erkannt wird, und der daher auch auf die beiden andern Stäbe Wirkungen 
ausüben muss, denen ähnlich, welche der erstere Versuch gezeigt hat; allein 
die Rechnung ergibt, und die Erfahrung bestätigt, dass diese Wirkungen zu 
klein ausfallen, um merklich zu sein. Noch weniger könnte also der schwächste 
Stab unter den dreien merkliche Wirkungen dieser Art erzeugen.

\newpage
%==============================================================================
% SECTION: Beobachtungen in Copenhagen und Mailand
% Pages 549-550
%==============================================================================
\section*{Beobachtungen der Variationen der Magnetnadel in Copenhagen und Mailand am 5.\ und 6.\ November 1834}
\addcontentsline{toc}{section}{Beobachtungen in Copenhagen und Mailand}

\textit{\textsc{Schumacher}. Astronomische Nachrichten Nr.\ 276. 1835 März 21.}

Seit der Vollendung des hiesigen magnetischen Observatorium werden hier 
unter andern regelmässig an gewissen im Voraus bestimmten Tagen die Variatio-
nen der magnetischen Declination ununterbrochen in kurzen Zeitintervallen be-
obachtet, wozu anfangs dieselben Termine gewählt waren, welche Herr von 
\textsc{Humboldt} schon vor mehreren Jahren angeordnet hatte. Seit dem vorigen Früh-
jahr haben sich schon ziemlich viele Astronomen und Physiker in den Besitz von 
ähnlichen Apparaten gesetzt, wie der hiesige ist, den ich an einem andern Orte 
hinlänglich beschrieben habe, und nehmen an jenen verabredeten Beobachtungen 
Theil. Gleich die ersten auf diese Weise gewonnenen gleichzeitigen Beobachtun-
gen am 4.\ und 5.\ Mai, in Göttingen und Waltershausen (einem Gute in der 
Gegend von Schweinfurt, wo Herr \textsc{Sartorius} mit einem zwar kleinen, aber sonst 
dem hiesigen ganz ähnlichen Apparat beobachtete), zeigten eine überaus merk-
würdige Harmonie in dem vielfach hin und her springenden Gange der Variatio-
nen, nicht blos in den grössern sondern auch in den geringern. Ähnliche Er-
folge haben sich seitdem in den spätern Terminen, wo Leipzig, Berlin, Braun-
schweig und Copenhagen Theil genommen haben, schon vielfach wiederholt; ei-
nige Proben sind in graphischen Darstellungen in \textsc{Poggendorff}'s Annalen mit-
getheilt.

So wie sich die Theilnahme an diesen verabredeten Beobachtungen immer 
weiter verbreiten wird, stehen natürlich immer interessantere und fruchtbarere 
Resultate zu erwarten. Ich wiederhole daher hier die bereits anderwärts gemachte 
Anzeige, dass wir, seitdem die Nothwendigkeit, in sehr kurzen Zeitintervallen 
zu beobachten, sich so klar herausgestellt hat, uns veranlasst gefunden haben, 
mit den Terminen eine Abänderung zu treffen, indem wir die Anzahl der Ter-
mine von 8 auf 6 im Jahr, und ihre Dauer von 44 auf 24 Stunden herabgesetzt 
haben. Die gegenwärtige Bestimmung ist der letzte Sonnabend jedes ungeraden 
Monats, vom Göttinger Mittag an bis zum Mittag des folgenden Tages. Es kom-
men zu diesen Hauptterminen noch jedesmal zwei Nebentermine, nemlich am 
nächstfolgenden Dienstag und Mittwoch Abends von 8 bis 10 Uhr. Umständli-
chere Nachricht, auch über Beobachtungsweise, findet man in \textsc{Poggendorff}'s An-
nalen Bd.\ 33. [S.\ 528 d.\ B.]

Das merkwürdigste bisher erhaltene Resultat bieten die gleichzeitigen Be-
obachtungen von Copenhagen und Mailand dar, am 5.\ und 6.\ November d.\ J., 
einem Termine nach dem frühern Arrangement, von dessen Abänderung die Be-
obachter an jenen Orten die Nachricht noch nicht erhalten hatten. In Copenha-
gen, wo jetzt unter Leitung des Herrn Etatsrath \textsc{Oersted} ein dem hiesigen ganz 
ähnliches magnetisches Observatorium errichtet ist, wurde eine Nadel von der-
selben Stärke, wie die hiesige, gebraucht (vier Pfund schwer); in Mailand be-
obachteten auf der dortigen Sternwarte die Herrn \textsc{Sartorius} und Doctor \textsc{Listing}, 
unter Beistand des Herrn \textsc{Kreil}, Eleven der Sternwarte, mit der schon oben er-
wähnten kleinern Nadel. Ich gestehe, dass ich, auch nach den vielen schon 
früher vorgekommenen Erfahrungen ähnlicher Art, doch durch die Grösse der 
Übereinstimmung an zwei mehr als 150 Meilen von einander entfernten Orten 
überrascht wurde. Der blosse Anblick der beigefügten graphischen Darstellung 
spricht hier für sich. Ich begleite dieselbe nur mit einigen Erläuterungen und 
Bemerkungen.

Da mir anfangs der Werth der Scalentheile in Copenhagen noch unbekannt 
war, so entwarf ich die Zeichnung nach solchem Maassstabe, dass die Anoma-
lien ungefähr gleich gross erscheinen, was ich erhielt, indem ich der Seite der 
Netzquadrate neun Scalentheile der Copenhagener, und drei der Mailänder Be-
obachtungen entsprechen liess. Ein Scalentheil in Copenhagen beträgt übrigens 
$21{,}''576$, einer in Mailand $29{,}''341$. Im Bogentheilen waren also die Copenhagener 
Bewegungen etwa 2,2 mal grösser als die Mailänder. Will man hieraus auf 
das Verhältniss der dabei thätigen Kräfte schliessen, so muss man nicht überse-
hen, dass diese Erscheinungen nur als Störungen der horizontalen erdmagneti-
schen Kraft an beiden Orten zu betrachten sind, und dass an einem Orte, wo 
letztere kleiner ist, eine gleiche störende Kraft grössere Änderungen hervorbrin-
gen muss, als an einem andern, wo jene grösser ist. Das Verhältniss der hori-
zontalen Kraft des Erdmagnetismus in Copenhagen und Mailand schätzte ich nach 
\textsc{Hansteen}'s schöner Karte im 7.\ Bande der Astronomischen Nachrichten\footnote{Im 9.\ Bande der Astronomischen Nachrichten hat dieser hochverdiente Naturforscher uns auch mit} 
wie

\end{document}
